\subsection{\texorpdfstring{复形 \(\mathsf{K}(u),\mathsf{K}.(u,\mathsf{C}),\mathsf{K}^{\prime}(u,\mathsf{C})\)}{复形 \textbackslash mathsf\{K\}(u),\textbackslash mathsf\{K\}.(u,\textbackslash mathsf\{C\}),\textbackslash mathsf\{K\}\^{}\{\textbackslash prime\}(u,\textbackslash mathsf\{C\})}}

设 A 为环，L 为 A-模， \(u: \mathbf{L} \to \mathbf{A}\) 为线性形式，且 \(\Lambda(\mathbf{L})\) 为 A-模 L 的外代数。对于 \(x \in \Lambda(\mathbf{L})\) ，令 \(d_u(x)\) 表示内积 \(x \sqsubseteq u\) （III，第 161 页，例）。根据前引文第 162 页公式 (60)，我们有

\[
d _ {u} \left(e _ {1} \wedge \dots \wedge e _ {n}\right) = \sum_ {i = 1} ^ {n} (- 1) ^ {i + 1} u \left(e _ {i}\right) e _ {1} \wedge \dots \wedge e _ {i - 1} \wedge e _ {i + 1} \wedge \dots \wedge e _ {n}\tag{1}
\]

对 L 中的 \(e_1, \ldots, e_n\) 成立。根据 III，第 164 及 165 页，映射 \(d_u: \symbf{\Lambda}(\mathbf{L}) \to \symbf{\Lambda}(\mathbf{L})\) 是次数为 \((-1)\) 且平方为零的反导子。它是 A-代数 \(\symbf{\Lambda}(\mathbf{L})\) 中扩张 \(u: \symbf{\Lambda}^1(\mathbf{L}) \to \symbf{\Lambda}^0(\mathbf{L})\) 的唯一反导子。

定义 1.--- 复形 \((\symbf{\Lambda}(\mathbf{L}), d_u)\) 记作 \(\mathbf{K}^{\mathbf{A}}(u)\) 或 \(\mathbf{K}(u)\) 。

注意 \(\mathbf{K}_{n}(u)=\symbf{\Lambda}^{n}(\mathbf{L})=\mathbf{K}^{-n}(u)\) 。显然 \(\mathbf{K}(u)\) 右零，且 \(\mathrm{H}_{0}(\mathbf{K}(u))=\mathrm{Coker}(u)=\mathrm{A}/\mathfrak{q}\) ，其中 \(\mathfrak{q}\) 是 A 的理想 \(u(\mathbf{L})\) 。

对于 A-模的每个复形 C，我们设

\[
\mathbf {K} _ {\cdot} ^ {\mathrm{A}} (u, \mathrm{C}) = \mathrm{C} \otimes_ {\mathrm{A}} \mathbf {K} ^ {\mathrm{A}} (u), \quad \mathbf {K} _ {\cdot} ^ {\cdot} (u, \mathrm{C}) = \operatorname{Homgr} _ {\mathrm{A}} (\mathbf {K} ^ {\mathrm{A}} (u), \mathrm{C}).
\]

\[
\mathrm{H} ^ {\mathsf {A}} (u, \mathrm{C}) = \mathrm{H} (\mathrm{C} \otimes_ {\mathsf {A}} \mathbf {K} ^ {\mathsf {A}} (u)) \quad \mathrm{H} _ {\mathsf {A}} ^ {\bullet} (u, \mathrm{C}) = \mathrm{H} (\mathrm{Homgr} _ {\mathsf {A}} (\mathbf {K} ^ {\mathsf {A}} (u), \mathrm{C})) ,
\]

\[
\mathrm{H} _ {r} ^ {\mathrm{A}} (u, \mathrm{C}) = \mathrm{H} _ {r} (\mathrm{C} \otimes_ {\mathrm{A}} \mathbf {K} ^ {\mathrm{A}} (u)), \quad \mathrm{H} _ {\mathrm{A}} ^ {r} (u, \mathrm{C}) = \mathrm{H} ^ {r} (\operatorname{Homgr} _ {\mathrm{A}} (\mathbf {K} ^ {\mathrm{A}} (u), \mathrm{C})) .
\]

因此，我们有 A-模的典范同态（X，第 62 页及第 82 页）

\[
\gamma_ {0}: \mathrm{H} _ {0} (\mathbf {C}) \otimes_ {\mathbf {A}} \mathrm{A} / \mathfrak {q} \rightarrow \mathrm{H} _ {0} ^ {\mathbf {A}} (u, \mathbf {C}),
\]

\[
\lambda^ {0}: \mathrm{H} _ {\mathrm{A}} ^ {0} (u, \mathrm{C}) \rightarrow \operatorname{Hom} _ {\mathrm{A}} (\mathrm{A} / \mathfrak {q}, \mathrm{H} ^ {0} (\mathrm{C})).
\]

引理 1. --- 若复形 C 右零（相应地，左零），则 \(\mathbf{K}^{\mathbf{A}}(u,\mathbf{C})\) （相应地， \(\mathbf{K}_{\mathbf{A}}^{*}(u,c)\) ）右零（相应地，左零），并且 \(\gamma_{0}\) （相应地， \(\lambda^{0}\) ）是双射的。

这由 X，第 62 页，命题 1 及第 82 页，命题 1 得出。

命题 1. --- 设 \(x \in L\) ；记 \(\mathbf{R}_{x}: y \mapsto x \wedge y\) 为由 \(x\) 在代数 \(\boldsymbol{\Lambda}(\mathbf{L})\) 中左乘的映射。则有 \(d_{u} \circ \mathbf{R}_{x} + \mathbf{R}_{x} \circ d_{u} = u(x)\) . \(1_{\boldsymbol{\Lambda}(\mathbf{L})} = u(x)_{\boldsymbol{\Lambda}(\mathbf{L})}\) .

事实上， \((d_u \circ R_x + R_x \circ d_u)(y) = d_u(x \wedge y) + x \wedge d_u(y)\) ；由于 \(d_u\) 是一个导子， \(d_u(x \wedge y) + x \wedge d_u(y) = d_u(x) \wedge y = u(x).y\) .

推论 1. --- 若 u 是满射的，则 \(\mathbf{K}(u)\) 同伦于零（X，第 34 页）， \(\mathbf{K}_{\cdot}^{\mathrm{A}}(u,\mathrm{C})\) 和 \(\mathbf{K}_{\cdot}^{\cdot}(\mathbf{u},\mathbf{C})\) 对每个复形 C 也如此。

事实上，存在 \(x \in \mathbf{L}\) 使得 \(u(x) = 1\) 。于是 \(\mathbf{K}(u)\) 由命题 1 同伦于零，因而 \(\mathbf{K}^{\mathrm{A}}(u, \mathbf{C})\) （X，第 64 页，命题 3）与 \(\mathbf{K}_{\mathrm{A}}^{*}(u, \mathbf{C})\) （X，第 83 页，命题 3）也如此。

推论 2. --- 设 C 是一个复形，Ann(C) 为其零化子。则 q + Ann(C) 零化 H \(^{A}\) (u, C) 与 H \(^{*}_{A}\) (u, C)。

对每个 \(\lambda \in q\) ，位似 \(\lambda_{\mathbf{K}(u)}\) 由该命题同伦于零，因此 \(1_{C} \otimes \lambda_{\mathbf{K}(u)}\) 以及 Homgr \((\lambda_{\mathbf{K}(u)}, 1_{C})\) 依据 X，第 64 页，命题 3 及 X，第 83 页，命题 3 也同伦于零；由此可知 \(\lambda\) 零化 H.(u, C) 与 H'(u, C)。若 \(\lambda \in \text{Ann}(C)\) ，则 \(1_{\mathbf{K}(u)} \otimes \lambda_{C}\) 以及 Homgr \((1_{\mathbf{K}(u)}, \lambda_{C})\) 为零。

假设 L 是投射的（相应地， \(\mathbf{K}(u)\) 在次数 \(>0\) 时零调）。则复形 \(\symbf{\Lambda}(\mathbf{L})\) 是投射的（相应地，是 A/q 的一个分解），依据 III，第 87 页，推论 2；由 X，第 102 页（相应地，第 100 页），因此对每个 A-模 M，我们有同态

\begin{enumerate}
\def\labelenumi{(\arabic{enumi})}
\setcounter{enumi}{1}
\tightlist
\item
\end{enumerate}

\[
\mathrm{H} _ {r} ^ {\mathbf {A}} (u, \mathrm{M}) \rightarrow \operatorname{Tor} _ {r} ^ {\mathbf {A}} (\mathrm{A} / \mathrm{q}, \mathrm{M}), \quad \operatorname{Ext} _ {\mathbf {A}} ^ {r} (\mathrm{A} / \mathrm{q}, \mathrm{M}) \rightarrow \mathrm{H} _ {\mathbf {A}} ^ {r} (u, \mathrm{M})\tag{3}
\]

\[
\operatorname{Tor} _ {r} ^ {\mathrm{A}} (\mathrm{A} / \mathrm{q}, \mathrm{M}) \rightarrow \mathrm{H} _ {r} ^ {\mathrm{A}} (u, \mathrm{M}), \quad \mathrm{H} ^ {r} (u, \mathrm{M}) \rightarrow \operatorname{Ext} _ {\mathrm{A}} ^ {r} (\mathrm{A} / \mathrm{q}, \mathrm{M}).
\]

若 L 是投射的且 \(\mathbf{K}(u)\) 在次数 \(>0\) 时零调，则上述同态 (2) 和 (3) 是双射且互为逆（X，第 102 页，命题 1）。

命题 2. --- 设 \((\mathbf{L}_{i})_{i\in I}\) 为一族 A-模，其中集合 I 有限且全序。设 u 为 \(\bigoplus_{i\in I}\mathbf{L}_{i}\) 上的线性形式， \(u_{i}\) 为它在 \(\mathbf{L}_{i}\) 上的限制。

A-代数的典范同构（III，第 84 页）

\[
g: \bigotimes_{\substack{i\in I}}\symbf{\Lambda}(\mathrm{L}_{i})\to \symbf{\Lambda}(\bigoplus_{i\in I}\mathrm{L}_{i})
\]

是复形 \(\otimes \mathbf{K}(u_i)\) （X，第 63 页）到复形 \(\mathbf{K}(u)\) 上的同构。

事实上，根据 X，第 64 页，注记 4，复形 \(\otimes_{i\in I}\mathbf{K}(u_i)\) 的微分 D 是一个反导子；反导子 \(d_u\) 与 \(g\circ D\circ g^{-1}\) （ \(\boldsymbol{\Lambda}(\oplus L_i)\) 的）在 \(\oplus L_i\) 上与映射 \(x\mapsto u(x)\) . 1（从 \(\oplus L_i\) 到 \(\boldsymbol{\Lambda}(\oplus L_i)\) ）一致，因此它们相等（III，第 128 页，推论）。

设 C 和 C' 是两个 A-模复形。我们有（X，第 63 页及第 99 页）复形的典范同构

\[
\mathbf {C} \otimes_ {\mathbf {A}} \left(\mathbf {C} ^ {\prime} \otimes_ {\mathbf {A}} \mathbf {K} (u)\right)\rightarrow \left(\mathbf {C} \otimes_ {\mathbf {A}} \mathbf {C} ^ {\prime}\right) \otimes_ {\mathbf {A}} \mathbf {K} (u)
\]

\[
\operatorname{Homgr} _ {\mathbf {A}} \left(\mathrm{C} ^ {\prime}, \operatorname{Homgr} _ {\mathbf {A}} (\mathbf {K} (u), \mathrm{C})\right)\rightarrow \operatorname{Homgr} _ {\mathbf {A}} \left(\mathrm{C} ^ {\prime} \otimes_ {\mathbf {A}} \mathbf {K} (u), \mathrm{C}\right),
\]

即，同构

\begin{enumerate}
\def\labelenumi{(\arabic{enumi})}
\setcounter{enumi}{3}
\tightlist
\item
\end{enumerate}

\[
\mathbf {C} \otimes_ {\mathbf {A}} \mathbf {K} ^ {\mathrm{A}} (u, \mathbf {C} ^ {\prime}) \rightarrow \mathbf {K} ^ {\mathrm{A}} (u, \mathbf {C} \otimes_ {\mathbf {A}} \mathbf {C} ^ {\prime})\tag{5}
\]

\[
\operatorname{Homgr} _ {\mathbf {A}} \left(\mathrm{C} ^ {\prime}, \mathbb {K} _ {\mathbf {A}} ^ {*} (u, \mathrm{C})\right)\rightarrow \operatorname{Homgr} _ {\mathbf {A}} \left(\mathbb {K} _ {\mathbf {A}} ^ {*} (u, \mathrm{C} ^ {\prime}), \mathrm{C}\right).
\]

在 (4) 和 (5) 中，取 \(\mathbf{C}^{\prime} = \mathbf{K}(u^{\prime})\) ，其中 \(u^{\prime}: L^{\prime} \to A\) 是 A-模 \(L^{\prime}\) 上的线性形式，并注意到 \(\mathbf{K}^{\mathrm{A}}(u, \mathbf{K}(u^{\prime}))\) 按定义等于 \(\mathbf{K}(u^{\prime}) \otimes_{\mathrm{A}} \mathbf{K}(u)\) ，根据命题 2，它等同于 \(\mathbf{K}(u^{\prime} \oplus u)\) ，其中 \(u^{\prime} \oplus u: L^{\prime} \oplus L \to A\) 是线性形式 \((x^{\prime}, x) \mapsto u^{\prime}(x^{\prime}) + u(x)\) 。于是我们得到复形的同构

\begin{enumerate}
\def\labelenumi{(\arabic{enumi})}
\setcounter{enumi}{5}
\tightlist
\item
\end{enumerate}

\[
\mathbf {K} _ {\bullet} ^ {\mathrm{A}} (u ^ {\prime} \oplus u, \mathrm{C}) \rightarrow \mathbf {K} _ {\bullet} ^ {\mathrm{A}} (u, \mathbf {K} _ {\bullet} ^ {\mathrm{A}} (u ^ {\prime}, \mathrm{C}))\tag{7}
\]

\[
\mathbf {K} _ {\mathrm{A}} ^ {\bullet} (u ^ {\prime}, \mathbf {K} _ {\mathrm{A}} ^ {\bullet} (u, \mathrm{C})) \rightarrow \mathbf {K} _ {\mathrm{A}} ^ {\bullet} (u ^ {\prime} \oplus u, \mathrm{C}).
\]

通过取同调，我们推出A-模的同构

\[
\begin{array}{c c} \mathrm{H} _ {r} ^ {\mathbf {A}} (u ^ {\prime} \oplus u, \mathrm{C}) \to \mathrm{H} _ {r} ^ {\mathbf {A}} (u, \mathbf {K} _ {\bullet} ^ {\mathbf {A}} (u ^ {\prime}, \mathrm{C}))  , & r \in \mathbf {Z}  , \\ \mathrm{H} _ {\mathbf {A}} ^ {r} (u ^ {\prime}, \mathbf {K} _ {\mathbf {A}} ^ {\bullet} (u, \mathrm{C})) \to \mathrm{H} _ {\mathbf {A}} ^ {r} (u ^ {\prime} \oplus u, \mathrm{C})  , & r \in \mathbf {Z}  . \end{array}
\]

最后，注意到由代数 \(\Lambda(L)\) 中的乘法导出的同态

\[
m: \mathbf {K} ^ {\mathrm{A}} (u) \otimes_ {\mathrm{A}} \mathbf {K} ^ {\mathrm{A}} (u) \rightarrow \mathbf {K} ^ {\mathrm{A}} (u)
\]

是复形态射（因为 \(d_{u}\) 是反导子）。假设 L 是秩为 \(n\) 的自由模，并与复形态射 \(\mathbf{K}^{\mathrm{A}}(u)\to\symbf{\Lambda}^{n}\mathrm{L}(-n)\) （它在次数 \(n\) 处为恒同）复合，由此导出一个复形态射

\[
\chi : \mathbf {K} ^ {\mathrm{A}} (u) \otimes_ {\mathrm{A}} \mathbf {K} ^ {\mathrm{A}} (u) \rightarrow \boldsymbol {\Lambda} ^ {n} \mathrm{L} (- n);
\]

根据 X，第 99 页，命题 12，该态射典范地对应一个复形态射

\[
\varphi : \mathbf {K} ^ {\mathrm{A}} (u) \rightarrow \operatorname{Homgr} _ {\mathrm{A}} \left(\mathbf {K} ^ {\mathrm{A}} (u), \boldsymbol {\Lambda} ^ {n} \mathrm{L} (- n)\right)
\]

它是双射的（III，第 87 页，公式 (20)）。对每个复形 C，由此导出一个复合同构

\[
\begin{array}{r l} \mathbf {K} _ {\bullet} ^ {\mathrm{A}} (u, \mathrm{C}) & = \mathrm{C} \otimes_ {\mathrm{A}} \mathbf {K} ^ {\mathrm{A}} (u) \xrightarrow {1 \otimes_ {\Phi}} \mathrm{C} \otimes \operatorname{Homgr} _ {\mathrm{A}} \left(\mathbf {K} ^ {\mathrm{A}} (u), \boldsymbol {\Lambda} ^ {n} \mathrm{L} (- n)\right) \to \\ & \to \operatorname{Homgr} _ {\mathrm{A}} \left(\mathbf {K} ^ {\mathrm{A}} (u), \mathrm{C} \otimes_ {\mathrm{A}} \boldsymbol {\Lambda} ^ {n} \mathrm{L} (- n)\right) = \mathbf {K} _ {\mathrm{A}} ^ {\bullet} (u, \mathrm{C} \otimes_ {\mathrm{A}} \boldsymbol {\Lambda} ^ {n} \mathrm{L} (- n)). \end{array}
\]

通过取同调，因此有典范同构

\[
\mathrm{H} _ {r} ^ {\mathbf {A}} (u, \mathrm{C}) \rightarrow \mathrm{H} _ {\mathbf {A}} ^ {n - r} (u, \mathrm{C} \otimes_ {\mathbf {A}} \boldsymbol {\Lambda} ^ {n} \mathrm{L}), \quad r \in \mathbf {Z}.\tag{8}
\]

注记。--- 1) * 当 L 为秩 n 的投射模时，前述结论仍然成立。~*

\begin{enumerate}
\def\labelenumi{\arabic{enumi})}
\setcounter{enumi}{1}
\tightlist
\item
  由于 L 是秩为 n 的自由模， \(\Lambda^{n}\) L 同构于 A，故有非典范同构 \(\mathrm{H}_{r}^{\mathrm{A}}(u,\mathrm{C})\to\mathrm{H}_{\mathrm{A}}^{n-r}(u,\mathrm{C})\) 。
\end{enumerate}

\subsection{函子性}

设 \(f: \mathbf{C} \to \mathbf{C}'\) 是复形的态射。我们用

\[
\begin{array}{l} \mathbf {K} _ {\bullet} ^ {\mathrm{A}} (u, f): \mathbf {K} _ {\bullet} ^ {\mathrm{A}} (u, \mathrm{C}) \to \mathbf {K} _ {\bullet} ^ {\mathrm{A}} (u, \mathrm{C} ^ {\prime}), \\ \mathbf {K} _ {\mathrm{A}} ^ {\bullet} (u, f): \mathbf {K} _ {\mathrm{A}} ^ {\bullet} (u, \mathrm{C}) \to \mathbf {K} _ {\mathrm{A}} ^ {\bullet} (u, \mathrm{C} ^ {\prime}), \end{array}
\]

表示复形的态射 \(f \otimes 1_{\mathbf{K}(u)}\) 和 \(\operatorname{Homgr}_{\mathbf{A}}(1_{\mathbf{K}(u)}, f)\) 。

我们用 \(\mathrm{H}_{\bullet}^{\mathrm{A}}(u,f):\mathrm{H}_{\bullet}^{\mathrm{A}}(u,\mathbf{C})\to\mathrm{H}_{\bullet}^{\mathrm{A}}(u,\mathbf{C}^{\prime}),\quad\mathrm{H}_{\mathrm{A}}^{\bullet}(u,f):\mathrm{H}_{\mathrm{A}}^{\bullet}(u,\mathbf{C})\to\mathrm{H}_{\mathrm{A}}^{\bullet}(u,\mathbf{C}^{\prime}),\quad\mathrm{H}_{r}^{\mathrm{A}}(u,f):\mathrm{H}_{r}^{\mathrm{A}}(u,\mathbf{C})\to\mathrm{H}_{r}^{\mathrm{A}}(u,\mathbf{C}^{\prime}),\mathrm{H}_{\mathrm{A}}^{r}(u,f):\mathrm{H}_{\mathrm{A}}^{r}(u,\mathbf{C})\to\mathrm{H}_{\mathrm{A}}^{r}(u,\mathbf{C}^{\prime})\) 表示在同调中诱导的态射。映射 \(f\mapsto\mathbf{K}_{\bullet}^{\mathrm{A}}(u,f)\) 是线性的；如果 \(g:\mathbf{C}^{\prime}\to\mathbf{C}^{\prime\prime}\) 是复形的另一个态射，我们有 \(\mathbf{K}_{\bullet}^{\mathrm{A}}(u,g\circ f)=\mathbf{K}_{\bullet}^{\mathrm{A}}(u,g)\circ\mathbf{K}_{\bullet}^{\mathrm{A}}(u,f)\) ；类似地，对于 \(\mathbf{K}_{\mathrm{A}}^{\bullet},\mathrm{H}_{\bullet}^{\mathrm{A}},\mathrm{H}_{\mathrm{A}}^{\bullet},\mathrm{H}_{r}^{\mathrm{A}},\mathrm{H}_{\mathrm{A}}^{r}\) .

设 \(0 \to \mathbf{C}' \xrightarrow{f} \mathbf{C} \xrightarrow{g} \dot{\mathbf{C}}'' \to 0\) 是复形的一个正合列。

\begin{enumerate}
\def\labelenumi{\alph{enumi})}
\tightlist
\item
  假设 L 是平坦的；则 \(\Lambda(L)\) 是平坦的（X，第 15 页，推论）。序列
\end{enumerate}

\[
0 \rightarrow \mathbf {K} _ {\bullet} ^ {\mathrm{A}} (u, C ^ {\prime}) \xrightarrow {\mathbf {K} _ {\bullet} ^ {\mathrm{A}} (u , f)} \mathbf {K} _ {\bullet} ^ {\mathrm{A}} (u, C) \xrightarrow {\mathbf {K} _ {\bullet} ^ {\mathrm{A}} (u , g)} \mathbf {K} _ {\bullet} ^ {\mathrm{A}} (u, C ^ {\prime \prime}) \rightarrow 0
\]

于是它是正合的，并（X，第 30 页）导出一个同调正合列。

\[
\dots \to \mathrm{H} _ {n} ^ {\mathbf {A}} (u, \mathrm{C} ^ {\prime}) \xrightarrow {\mathrm{H} _ {n} (u , f)} \mathrm{H} _ {n} ^ {\mathbf {A}} (u, \mathrm{C}) \xrightarrow {\mathrm{H} _ {n} (u , g)} \mathrm{H} _ {n} ^ {\mathbf {A}} (u, \mathrm{C} ^ {\prime \prime}) \xrightarrow {\hat {\sigma} _ {n}} \mathrm{H} _ {n - 1} ^ {\mathbf {A}} (u, \mathrm{C} ^ {\prime}) \to \dots .
\]

\begin{enumerate}
\def\labelenumi{\alph{enumi})}
\setcounter{enumi}{1}
\tightlist
\item
  假设 L 是投射的；则 \(\Lambda(L)\) 是投射的。序列
\end{enumerate}

\[
0 \rightarrow \mathbf {K} _ {\mathrm{A}} ^ {*} (u, \mathrm{C} ^ {\prime}) \xrightarrow {\mathbf {K} _ {\mathrm{A}} ^ {*} (u , f)} \mathbf {K} _ {\mathrm{A}} ^ {*} (u, \mathrm{C}) \xrightarrow {\mathbf {K} _ {\mathrm{A}} ^ {*} (u , g)} \mathbf {K} _ {\mathrm{A}} ^ {*} (u, \mathrm{C} ^ {\prime \prime}) \rightarrow 0
\]

于是它是正合的，并导出一个同调正合列。

\[
\dots \rightarrow \mathrm{H} _ {\mathrm{A}} ^ {n} (u, \mathrm{C} ^ {\prime}) \xrightarrow {\mathrm{H} ^ {n} (u , f)} \mathrm{H} _ {\mathrm{A}} ^ {n} (u, \mathrm{C}) \xrightarrow {\mathrm{H} ^ {n} (u , g)} \mathrm{H} _ {\mathrm{A}} ^ {n} (u, \mathrm{C} ^ {\prime \prime}) \xrightarrow {\partial^ {n}} \mathrm{H} _ {\mathrm{A}} ^ {n + 1} (u, \mathrm{C} ^ {\prime}) \rightarrow \dots .
\]

设 \(\rho: A \to A'\) 是一个环同态， \(L'\) 为 \(A'\) -模 \(A' \otimes_A L\) ， \(u': L' \to A'\) 为线性形式 \(1 \otimes u\) 。典范双射同态（III，第 83 页，命题 8）

\[
\psi : \Lambda_ {A ^ {\prime}} (A ^ {\prime} \otimes_ {A} L) \rightarrow A ^ {\prime} \otimes_ {A} \Lambda_ {A} (L)
\]

是 A'-模复形的同构。我们推出：

\begin{enumerate}
\def\labelenumi{\arabic{enumi})}
\tightlist
\item
  对每个 A'-模复形 C'，一个 A-模复形的同构
\end{enumerate}

\[
\mathbf {K} _ {\bullet} ^ {\mathrm{A} ^ {\prime}} \left(u ^ {\prime}, \mathrm{C} ^ {\prime}\right)\rightarrow \mathbf {K} _ {\bullet} ^ {\mathrm{A}} \left(u, \mathrm{C} ^ {\prime}\right),
\]

\begin{enumerate}
\def\labelenumi{(\arabic{enumi})}
\setcounter{enumi}{8}
\tightlist
\item
\end{enumerate}

它由图表

\[
\mathrm{C} ^ {\prime} \otimes_ {\mathrm{A} ^ {\prime}} \left(\boldsymbol {\Lambda} _ {\mathrm{A} ^ {\prime}} \left(\mathrm{A} ^ {\prime} \otimes_ {\mathrm{A}} \mathrm{L}\right)\right) \xrightarrow {1 _ {\mathrm{C} ^ {\prime}} \otimes \psi} \mathrm{C} ^ {\prime} \otimes_ {\mathrm{A} ^ {\prime}} \mathrm{A} ^ {\prime} \otimes_ {\mathrm{A}} \boldsymbol {\Lambda} _ {\mathrm{A}} (\mathrm{L}) \xrightarrow {\varphi} \mathrm{C} ^ {\prime} \otimes_ {\mathrm{A}} \boldsymbol {\Lambda} _ {\mathrm{A}} (\mathrm{L})
\]

复合而成，其中 φ 是典范双射（III，第 85 页，命题 14）；

\begin{enumerate}
\def\labelenumi{\arabic{enumi})}
\setcounter{enumi}{1}
\tightlist
\item
  对每个 A-模复形 C，一个 A'-模复形的同构
\end{enumerate}

\[
\mathbf {K} _ {\bullet} ^ {\mathrm{A}} (u, \mathrm{A} ^ {\prime} \otimes_ {\mathrm{A}} \mathrm{C}) \rightarrow \mathrm{A} ^ {\prime} \otimes_ {\mathrm{A}} \mathbf {K} _ {\bullet} ^ {\mathrm{A}} (u, \mathrm{C}),
\]

由此得到A'-模的同态

\[
\mathrm{A} ^ {\prime} \otimes_ {\mathrm{A}} \mathrm{H} _ {n} ^ {\mathrm{A}} (u, \mathrm{C}) \rightarrow \mathrm{H} _ {n} ^ {\mathrm{A} ^ {\prime}} (u ^ {\prime}, \mathrm{A} ^ {\prime} \otimes_ {\mathrm{A}} \mathrm{C}),
\]

当 A' 在 A 上平坦时，这些同态是双射的（X，第 66 页，推论 2）。

设 L' 为一个 A-模， \(u': L' \to A\) 为一个线性形式， \(f: L \to L'\) 为一个 A-同态，使得 \(u' \circ f = u\) 。由 III，第 161 页，公式 (55)，可知同态 \(\symbf{\Lambda}(f): \symbf{\Lambda}(\mathrm{L}) \to \symbf{\Lambda}(\mathrm{L}')\) 满足 \(d_u \circ \symbf{\Lambda}(f) = \symbf{\Lambda}(f) \circ d_{u'}\) ，从而定义了复形的态射 \(\symbf{\Lambda}(u): \mathbf{K}^{\mathbf{A}}(u) \to \mathbf{K}^{\mathbf{A}}(u')\) 。若 C 是一个 A-复形，我们推出复形的态射

\[
1 _ {C} \otimes \boldsymbol {\Lambda} (u): \mathbf {K} ^ {\mathrm{A}} (u, C) \rightarrow \mathbf {K} ^ {\mathrm{A}} \left(u ^ {\prime}, C\right) \quad \text { and }\quad \operatorname{Homgr} \left(\boldsymbol {\Lambda} (u), 1 _ {C}\right): \mathbf {K} _ {A} ^ {*} \left(u ^ {\prime}, C\right)\rightarrow \mathbf {K} _ {A} ^ {*} (u, C)
\]

若f是双射，则所有这些态射都是同构。

\subsection{\texorpdfstring{例1：复形 \(\mathbf{S}(\mathbf{L})\otimes_{\mathbf{A}}\boldsymbol{\Lambda}(\mathbf{L})\)}{例1：复形 \textbackslash mathbf\{S\}(\textbackslash mathbf\{L\})\textbackslash otimes\_\{\textbackslash mathbf\{A\}\}\textbackslash boldsymbol\{\textbackslash Lambda\}(\textbackslash mathbf\{L\})}}

设 A 为环，L 为 A-模， \(\mathbf{S}(\mathrm{L})\) 为其对称代数， \(\mathbf{S}(\mathrm{L}) \otimes_{\mathrm{A}} \mathrm{L}\) 为经纯量扩张得到的 \(\mathbf{S}(\mathrm{L})\) -模， \(u: \mathbf{S}(\mathrm{L}) \otimes_{\mathrm{A}} \mathrm{L} \to \mathbf{S}(\mathrm{L})\) 为线性形式，使得 \(u(s \otimes x) = sx\) 对 \(s \in \mathbf{S}(\mathrm{L})\) 、 \(x \in \mathrm{L}\) 成立。根据 \(\mathbf{S}(\mathrm{L})\) -模的典范同构（III，第 83 页，命题 8）

\[
\boldsymbol {\Lambda} _ {\mathbf {S} (L)} (\mathbf {S} (L) \otimes_ {A} L) \rightarrow \mathbf {S} (L) \otimes_ {A} \boldsymbol {\Lambda} (L),
\]

复形 \(\mathbf{K}^{\mathbf{S}(\mathrm{L})}(u)\) 的微分被搬运到映射

\[
d: \mathbf {S} (\mathrm{L}) \otimes_ {\mathrm{A}} \boldsymbol {\Lambda} (\mathrm{L}) \rightarrow \mathbf {S} (\mathrm{L}) \otimes_ {\mathrm{A}} \boldsymbol {\Lambda} (\mathrm{L})
\]

使得对 L 中的 \(x_{1},\ldots,x_{p},y_{1},\ldots,y_{q}\) ，我们有

\[
\begin{array}{l} d \left(\left(x _ {1} \dots x _ {p}\right) \otimes \left(y _ {1} \wedge \dots \wedge y _ {q}\right)\right) \\ = \sum_ {i = 1} ^ {q} (- 1) ^ {i + 1} y _ {i} x _ {1} \dots x _ {p} \otimes \left(y _ {1} \wedge \dots \wedge y _ {i - 1} \wedge y _ {i + 1} \wedge \dots \wedge y _ {q}\right). \end{array}
\]

注意 \(d\) 把 \(\mathbf{S}^{p}(\mathrm{L})\otimes\symbf{\Lambda}^{q}(\mathrm{L})\) 映到 \(\mathbf{S}^{p+1}(\mathrm{L})\otimes\symbf{\Lambda}^{q-1}(\mathrm{L})\) ，因此 A-模复形 \(\mathbf{S}(\mathrm{L})\otimes\symbf{\Lambda}(\mathrm{L})\) 分解为下述图表所描述的复形的直和：

\[
\left(\mathcal {E} _ {n}\right): 0 \rightarrow \mathbf {S} ^ {0} L \otimes_ {\mathrm{A}} \boldsymbol {\Lambda} ^ {n} L \rightarrow \mathbf {S} ^ {1} L \otimes_ {\mathrm{A}} \boldsymbol {\Lambda} ^ {n - 1} L \rightarrow \dots \rightarrow \mathbf {S} ^ {n} L \otimes_ {\mathrm{A}} \boldsymbol {\Lambda} ^ {0} L \rightarrow 0, n \in \mathbf {N}.
\]

若 A-模 L 是有限族 \((\mathbf{L}_{i})_{i\in I}\) 的直和，其中 I 是全序的，则典范双射

\[
\bigotimes_ {i \in I} ^ {\varepsilon} \left(\mathbf {S} (L _ {i}) \otimes_ {A} \boldsymbol {\Lambda} (L _ {i})\right) \to \mathbf {S} (L) \otimes_ {A} \boldsymbol {\Lambda} (L)
\]

是 A-模复形的同构（这由 X，第 148 页，命题 2 或上面的公式 (9) 推出）。

命题 3. --- 若 A-模 L 是平坦的，则上述序列 \((\mathcal{E}_n)\) 对 \(n > 0\) 是正合的。

\begin{enumerate}
\def\labelenumi{\alph{enumi})}
\tightlist
\item
  首先注意，若 \(p_{L}\) 是复合同态
\end{enumerate}

\[
\mathbf {S} (\mathrm{L}) \otimes \boldsymbol {\Lambda} (\mathrm{L}) \xrightarrow {\alpha} \mathbf {S} ^ {0} (\mathrm{L}) \otimes \boldsymbol {\Lambda} ^ {0} (\mathrm{L}) \xrightarrow {\beta} \mathrm{A},
\]

其中 \(\alpha\) 是诸典范投影的张量积， \(\beta\) 是典范同构；只需证明 \(\mathbf{H}(p_{\mathrm{L}})\) 是双射的。

\begin{enumerate}
\def\labelenumi{\alph{enumi})}
\setcounter{enumi}{1}
\item
  若L=0或L=A，命题显然成立。
\item
  假设 L 是有限秩的自由模；将它写成秩 1 的自由 A-模的直和 \(L_{1} \oplus \ldots \oplus L_{n}\) 。根据命题前的注记，复形 \(\mathbf{S}(\mathrm{L}) \otimes \symbf{\Lambda}(\mathrm{L})\) 同构于 n 个自由复形 \(\mathbf{S}(\mathrm{L}_{i}) \otimes \symbf{\Lambda}(\mathrm{L}_{i})\) 的张量积，后者的同调由 b) 知是自由的。由 X，第 79 页，推论 4，典范同态
\end{enumerate}

\[
\gamma : \bigoplus_ {i = 1} ^ {n} \mathrm{H} (\mathbf {S} (\mathrm{L} _ {i}) \otimes \boldsymbol {\Lambda} (\mathrm{L} _ {i})) \rightarrow \mathrm{H} (\mathbf {S} (\mathrm{L}) \otimes \boldsymbol {\Lambda} (\mathrm{L}))
\]

是双射的。由 b) 知 \(\mathrm{H}(p_{\mathrm{L}_i})\) 对所有 \(i\) 都是双射的。由于 \(\bigotimes_{i=1}^{n} \mathrm{H}(p_{\mathrm{L}_i}) = \mathrm{H}(p_{\mathrm{L}}) \circ \gamma\) ，

\[
\mathrm{H} (p _ {\mathrm{L}}) \text {   is   bijective }  .
\]

\begin{enumerate}
\def\labelenumi{\alph{enumi})}
\setcounter{enumi}{3}
\tightlist
\item
  在一般情况下，L 是由有限秩自由模构成的滤过归纳系 \((\mathbf{L}_{i})_{i\in I}\) 的归纳极限（X，第 14 页，定理 1）。由于典范双射同态
\end{enumerate}

\[
\varinjlim \mathbf {S} (\mathrm{L} _ {i}) \otimes \boldsymbol {\Lambda} (\mathrm{L} _ {i}) \rightarrow \mathbf {S} (\mathrm{L}) \otimes \boldsymbol {\Lambda} (\mathrm{L})
\]

是复形的同构，命题由 X，第 28 页，命题 1 得出。

注记。--- 1) 我们将在下文（X，第 158 页，例）看到上述 c) 部分的另一个证明。

\begin{enumerate}
\def\labelenumi{\arabic{enumi})}
\setcounter{enumi}{1}
\tightlist
\item
  如果 A 是一个 Q-代数，那么命题 3 的结论在不对 L 作任何假设的情况下仍然成立（参见 X，第 206 页，习题 1）。
\end{enumerate}

设 G 是一个群， \(\rho: G \to \mathbf{GL}(L)\) 是 G 在平坦 A-模 L 中的一个线性表示。则 \((\mathcal{E}_n)\) 是线性表示的正合列。假设 L 是有限生成投射模，并记 \(R_A(G)\) 为 G 在有限生成投射 A-模中的表示环。由命题 3 可知，我们在 \(R_{A}(G)\) 中有关系

\[
\sum_ {i = 0} ^ {n} (- 1) ^ {i} \left[ \mathbf {S} ^ {i} (\mathrm{L}) \right] \left[ \boldsymbol {\Lambda} ^ {n - i} (\mathrm{L}) \right] = 0, \quad n > 0.\tag{10}
\]

如果我们考虑形式幂级数

\[
s (\mathrm{T}) = \sum_ {i = 0} ^ {\infty} \left[ \mathbf {S} ^ {i} (\mathrm{L}) \right] \mathrm{T} ^ {i} \in \mathrm{R} _ {\mathbf {A}} (\mathrm{G}) [ [ \mathrm{T} ] ],
\]

\[
\lambda (\mathrm{T}) = \sum_ {i = 0} ^ {\infty} \left[ \boldsymbol {\Lambda} ^ {i} (\mathrm{L}) \right] \mathrm{T} ^ {i} \in \mathrm{R} _ {\mathbf {A}} (\mathrm{G}) [ [ \mathrm{T} ] ],
\]

关系式 (10) 可写作

\begin{enumerate}
\def\labelenumi{(\arabic{enumi})}
\setcounter{enumi}{10}
\tightlist
\item
\end{enumerate}

\[
s (\mathrm{T}) \lambda (- \mathrm{T}) = 1 _ {*}.
\]

\subsection{例2：自由模的情形}

设 \(k\) 为一个环， \(\mathbf{M}\) 为 \(k\) -模， \(\mathbf{I}\) 为一个集合， \(p\) 为一个整数 \(\geqslant 0\) 。称映射 \(\boldsymbol{m}:1^{p}\to\mathbf{M}\) 是交错的，如果它满足以下两个条件：

\begin{enumerate}
\def\labelenumi{\alph{enumi})}
\tightlist
\item
  对每个置换 \(\sigma \in \mathfrak{S}_p\) 以及每个序列 \((\alpha_1, \ldots, \alpha_p) \in \mathrm{I}^p\) ，我们有
\end{enumerate}

\[
\boldsymbol {m} \left(\alpha_ {\sigma (1)}, \dots , \alpha_ {\sigma (p)}\right) = \varepsilon_ {\sigma} \boldsymbol {m} \left(\alpha_ {1}, \dots , \alpha_ {p}\right),
\]

\begin{enumerate}
\def\labelenumi{\alph{enumi})}
\setcounter{enumi}{1}
\tightlist
\item
  对每个序列 \((\alpha_{1},\ldots,\alpha_{p})\in\mathbb{I}^{p}\) ，若指标 \(\alpha_{1},\ldots,\alpha_{p}\) 中有两个相等，则 \(m(\alpha_{1},\ldots,\alpha_{p})=0\) 。
\end{enumerate}

（在 I 是 k-模且 m 是多线性的情形下，就恢复了 III，第 80 页中引入的概念。）

假设 I 是有限的，并用 \(C_{I}^{p}(M)\) 表示从 \(I^{p}\) 到 M 的交错映射的 k-模。

设 \(L_{0}\) 是一个 k-模， \((e_{i})_{i\in I}\) 为属于 \(L_{0}\) 的一个元素族；我们定义两个 k-线性映射

\[
\begin{array}{l} g: \operatorname{Hom} _ {k} (\symbf {\Lambda} ^ {p} \mathrm{L} _ {0}, \mathrm{M}) \to \mathrm{C} _ {\mathrm{I}} ^ {p} (\mathrm{M}) \\ h: \mathrm{C} _ {\mathrm{I}} ^ {p} (\mathrm{M}) \to \mathrm{M} \otimes_ {k} \symbf {\Lambda} ^ {p} \mathrm{L} _ {0} \end{array}
\]

如下：如果 \(f \in \operatorname{Hom}_k(\symbf{\Lambda}^p \mathbf{L}_0, \mathbf{M})\) ，我们设

\[
g (f) (\alpha_ {1}, \dots , \alpha_ {p}) = f (e _ {\alpha_ {1}} \wedge \dots \wedge e _ {\alpha_ {p}});
\]

设 \(m \in C_{I}^{p}(M)\) ，我们定义 \(h(m) \in M \otimes_{k} \Lambda^{p} L_{0}\) 。对每个元素 \((\alpha_{1}, \ldots, \alpha_{p})\)

属于 \(I^{p}\) ，元素 \(m(\alpha_{1},\ldots,\alpha_{p})\otimes(e_{\alpha_{1}}\wedge\ldots\wedge e_{\alpha_{p}})\) （属于 \(\Lambda^{p}L_{0}\otimes_{k}M\) ）当 Card \(\{\alpha_{1},\ldots,\alpha_{p}\}<p\) 时为零，并且与指标 \(\alpha_{1},\ldots,\alpha_{p}\) 的顺序无关，如果

\[
\text { Card } \{\alpha_ {1}, \dots , \alpha_ {p} \} = p.
\]

的顺序无关。它只依赖于 I 的子集 \(\mathbf{J} = (\alpha_{1},\ldots ,\alpha_{p})\) ；将其记作 \(h_\mathrm{J}(\pmb {m})\) ；我们有 \(h_\mathrm{J}(\pmb {m}) = 0\) ，如果 Card (J) \(< p\) ；然后我们设：

\[
h (\boldsymbol {m}) = \sum_ {\mathrm{J}} h _ {\mathrm{J}} (\boldsymbol {m}),
\]

其中J遍历I的所有含p个元素的子集。

引理 2. --- 若 k-模 \(\mathbf{L}_{0}\) 以 \((e_{i})_{i\in I}\) 为基，则 \(k\) -线性映射 \(g\) 与 \(h\) 都是双射的。

这由 III，第 79 页，定理 1 得出。

现在设 M 为一个 k-模，且 \(\boldsymbol{x} = (x_{i})_{i \in I}\) 为 M 的一族两两交换的 k-自同态。考虑多项式环 \(\mathbf{A} = k[(\mathbf{X}_{i})_{i \in I}]\) 并赋予 M 以 A-模结构，使得 \(\mathrm{P}(\mathbf{X}_{i}) m = \mathrm{P}(x_{i}) m\) 对 \(P \in A\) 和 \(m \in M\) 成立。另一方面，设 L 为自由 A-模 \(A^{I}, (e_{i})_{i \in I}\) 为其典范基， \(u : L \to A\) 为把 \(e_{i}\) 映到 \(X_{i}\) （对所有 \(i \in I\) ）的线性形式。考虑 k-模复形 \(\mathbf{K}^{\mathbf{A}}(u, \mathbf{M})\) 与 \(\mathbf{K}_{\mathbf{A}}^{\bullet}(u, \mathbf{M})\) ；我们有典范同构

\[
\begin{array}{c} \mathbf {K} _ {\mathbf {A}} ^ {p} (u, \mathrm{M}) = \operatorname{Hom} _ {\mathbf {A}} \left(\boldsymbol {\Lambda} _ {\mathbf {A}} ^ {p} (\mathrm{A} ^ {\mathrm{I}}), \mathrm{M}\right) \to \operatorname{Hom} _ {k} \left(\boldsymbol {\Lambda} _ {k} ^ {p} (k ^ {\mathrm{I}}), \mathrm{M}\right), \\ \mathrm{M} \otimes_ {k} \boldsymbol {\Lambda} _ {k} ^ {p} (k ^ {\mathrm{I}}) \to \mathrm{M} \otimes_ {\mathbf {A}} \boldsymbol {\Lambda} _ {\mathbf {A}} ^ {p} (\mathrm{A} ^ {\mathrm{I}}) = \mathbf {K} _ {p} ^ {\mathbf {A}} (u, \mathrm{M}); \end{array}
\]

由此，通过与同构 \(g\) 和 \(h\) 复合，得到 \(k\) -模的同构

\[
\begin{array}{l} \theta^ {p}: \mathbf {K} _ {\mathrm{A}} ^ {p} (u, \mathrm{M}) \to \mathrm{C} _ {\mathrm{I}} ^ {p} (\mathrm{M}), \\ \theta_ {p}: \mathrm{C} _ {\mathrm{I}} ^ {p} (\mathrm{M}) \to \mathbf {K} _ {p} ^ {\mathrm{A}} (u, \mathrm{M}). \end{array}
\]

我们用

\[
\begin{array}{l} \hat {c} ^ {p}: \mathbf {C} _ {\mathrm{I}} ^ {p} (\mathbf {M}) \to \mathbf {C} _ {\mathrm{I}} ^ {p + 1} (\mathbf {M}), \\ \hat {c} _ {p}: \mathbf {C} _ {\mathrm{I}} ^ {p} (\mathbf {M}) \to \mathbf {C} _ {\mathrm{I}} ^ {p - 1} (\mathbf {M}), \end{array}
\]

的 \(k\) -同态，它们是通过搬运 \(\mathbf{K}_{\mathbf{A}}^{\bullet}(u, \mathbf{M})\) 与 \(\mathbf{K}_{\mathbf{A}}^{\mathbf{A}}(u, \mathbf{M})\) 的微分经同构 \(\theta\) 而得到的。例如，我们有：

\[
\left(\partial^ {p} \boldsymbol {m}\right) \left(\alpha_ {1}, \dots , \alpha_ {p + 1}\right) = \sum_ {j = 1} ^ {p + 1} (- 1) ^ {j + 1} x _ {\alpha_ {j}} \boldsymbol {m} \left(\alpha_ {1}, \dots , \alpha_ {j - 1}, \alpha_ {j + 1}, \dots , \alpha_ {p + 1}\right).\tag{12}
\]

由 \(C_{I}^{p}(M)\) 与 \(\partial^{p}\) （相应地， \(\partial_{p}\) ）构成的复形记作 \(\mathbf{K}^{\prime}(\mathbf{x},\mathbf{M})\) （相应地， \(\mathbf{K}(\mathbf{x},\mathbf{M})\) ），称为与模 M 及自同态序列

\((x_{1}, \ldots, x_{n})\) 相伴的上升（相应地，下降）科斯居尔（Koszul）复形。因此我们有 k-模复形的同构

\[
\begin{array}{l} \theta^ {\bullet}: \mathbf {K} _ {\mathbf {A}} ^ {\bullet} (u, \mathbf {M}) \to \mathbf {K} ^ {\bullet} (x, \mathbf {M}), \\ \theta_ {\bullet}: \mathbf {K}. (x, \mathbf {M}) \to \mathbf {K} _ {\bullet} ^ {\mathbf {A}} (u, \mathbf {M}). \end{array}
\]

注记。--- 反过来，设 B 是一个 \(k\) -代数，L 是一个以 \((e_i)_{i \in I}\) 为基的自由 B-模，M 是一个 B-模。线性形式 \(u: L \to B\) 的给定等价于 B 中元素的一族 \(x = (x_i)_{i \in 1}\) ，对应关系为 \(x_i = u(e_i)\) 。作为 \(k\) -模复形， \(\mathbf{K}_{\mathrm{B}}^{\bullet}(u, \mathrm{M})\) （相应地， \(\mathbf{K}^{\mathrm{B}}(u, \mathrm{M})\) ）通过同构 \(\theta^{\bullet}\) （相应地， \(\theta\) .）与科斯居尔复形 \(\mathbf{K}'(x, \mathrm{M})\) （相应地， \(\mathbf{K}^{\bullet}(x, \mathrm{M})\) ）等同。例如 \(\mathbf{K}^{\mathrm{B}}(u)\) 与 \(\mathbf{K}'(x, \mathrm{B})\) 等同。

我们如同第 1 节（X，第 147 页）那样引入记号 H.(x, M)、H.(x, M) 等，由于模 A \(^{I}\) 是自由模，第 1 节和第 2 节的所有结果均可经必要修改后适用。例如，我们有同构

\[
\begin{array}{l} \mathrm{H} _ {0} (\boldsymbol {x}, \mathrm{M}) \to \mathrm{M} / (\boldsymbol {x}) \mathrm{M} \\ \mathrm{H} ^ {0} (\boldsymbol {x}, \mathrm{M}) \to \mathrm{Hom} _ {\mathrm{A}} \left(\mathrm{A} / (\boldsymbol {x}), \mathrm{M}\right), \end{array}
\]

其中 \((\pmb{x})\) 表示 A 的理想 \(\sum \mathbf{A}x_{i}\) 。又例如，若 \(\mathbf{K}.(\pmb{x},\mathbf{A})\) 在次数 \(>0\) 时零调，我们有同构

\[
\begin{array}{l}\mathrm{H} _ {r} (\boldsymbol {x}, \mathrm{M}) \rightarrow \operatorname{Tor} _ {r} ^ {\mathrm{A}} (k, \mathrm{M}),\\\operatorname{Ext} _ {\mathrm{A}} ^ {r} (k, \mathrm{M}) \rightarrow \mathrm{H} ^ {r} (\boldsymbol {x}, \mathrm{M}).\end{array}
\]

最后，假设 I 是（有限的且）全序的，例如 \(I = \{1, \ldots, n\}\) ；让我们把 \(\symbf{\Lambda}^{n}(\mathbf{A}^{\mathrm{I}})\) 借助基元素 \(e_{1} \wedge \ldots \wedge e_{n}\) 等同于 A，并翻译 X，第 149 页的同构 \(\mathbf{K}_{p}^{\mathrm{A}}(u, \mathrm{M}) \to \mathbf{K}_{\mathrm{A}}^{n-p}(u, \mathrm{M})\) 。通过 \((\theta_{p})\) 与 \((\theta^{n-p})\) 的搬运，它变为同构

\[
\mathrm{C} _ {\mathrm{I}} ^ {p} (\mathbf {M}) \rightarrow \mathrm{C} _ {\mathrm{I}} ^ {n - p} (\mathbf {M})
\]

它把 \(m \in C_{I}^{p}(M)\) 对应到满足下列条件的元素 \(\check{m}\) （属于 \(C_{1}^{n-p}(M)\) ）：

\[
\boldsymbol {m} \left(\alpha_ {1}, \dots , \alpha_ {p}\right) = \check {\boldsymbol {m}} \left(\beta_ {1}, \dots , \beta_ {n - p}\right)\tag{13}
\]

若 \((\alpha_{1},\ldots,\alpha_{p},\beta_{1},\ldots,\beta_{n-p})\) 是 \(\{1,\ldots,n\}\) 的一个偶置换。我们也注意到，当 \(I=\{1,\ldots,n\}\) 时，可以把 \(C_{I}^{p}(M)\) 与元素取自 M 的族 \(m(\alpha_{1},\ldots,\alpha_{p})\) （其中 \(\alpha_{1}<\alpha_{2}<\ldots<\alpha_{p}\) ）的集合等同；公式 (12) 仍然成立，关系式 (13) 亦然。

例。--- 取 \(\mathbf{M} = k[\mathbf{T}_1, \ldots, \mathbf{T}_n]\) ；科斯居尔复形 \(\mathbf{K}\) （ \(\partial/\partial\mathbf{T}\) , \(\mathbf{M}\) ）——它与自同态序列（ \(\partial/\partial\mathbf{T}_1, \ldots, \partial/\partial\mathbf{T}_n\) ）相伴——等同于 \(k[x_{1},\ldots,x_{n}]\) 在 \(k\) 上的德拉姆复形（X，第 44 页）：对 \(\boldsymbol{m}\in\mathbf{C}_{\{1,\ldots,n\}}^{p}(\mathbf{M})\) ，对应以微分形式

\[
\omega (\boldsymbol {m}) = \sum_ {\alpha_ {1} <   \dots <   \alpha_ {p}} \boldsymbol {m} \left(\alpha_ {1}, \dots , \alpha_ {p}\right) d x _ {\alpha_ {1}} \wedge \dots \wedge d x _ {\alpha_ {p}},
\]

参见公式 (12) 及第 44 页例 1。

\subsection{例 3：情形 L = A}

若 L = A，令 \(u(1) = x \in \mathbf{A}\) 。复形 \(\mathbf{K}(u)\) 的长度为 1，且我们有 \(\mathbf{K}_{0}(u) = \mathbf{K}_{1}(u) = \mathbf{A}\) 和 \(d_{1}(a) = xa\) ，因此对每个 A-模 M， \(\mathbf{K}_{0}(u, \mathbf{M})\) , \(\mathbf{K}_{1}(u, \mathbf{M})\) , \(\mathbf{K}^{0}(u, \mathbf{M})\) 和 \(\mathbf{K}^{1}(u, \mathbf{M})\) 与 M 等同，微分

\[
d _ {1}: \mathbf {K} _ {1} (u, \mathrm{M}) \rightarrow \mathbf {K} _ {0} (u, \mathrm{M}) \quad \text { and }\quad d ^ {0}: \mathbf {K} ^ {0} (u, \mathrm{M}) \rightarrow \mathbf {K} ^ {1} (u, \mathrm{M})
\]

是 \(m \mapsto xm\) 。因此我们有同构

\[
\begin{array}{l} \mathrm{H} _ {0} (x, \mathrm{M}) \to \mathrm{A} / x \mathrm{A} \otimes_ {\mathrm{A}} \mathrm{M} \leftarrow \mathrm{H} ^ {1} (x, \mathrm{M}), \\ \mathrm{H} _ {1} (x, \mathrm{M}) \to \mathrm{Hom} _ {\mathrm{A}} (\mathrm{A} / x \mathrm{A}, \mathrm{M}) \leftarrow \mathrm{H} ^ {0} (x, \mathrm{M}). \end{array}
\]

引理3. --- 设 K 是一个复形，使得 \(\mathbf{K}_i = 0\) 当 \(i \neq 0, 1\) ，并设 C 是一个复形， \(p\) 为一个整数。

\begin{enumerate}
\def\labelenumi{\alph{enumi})}
\tightlist
\item
  若 \(K\) 是平坦的，则对所有 \(p \in \mathbf{N}\) ，我们有正合列
\end{enumerate}

\[
0 \rightarrow \mathrm{H} _ {0} (\mathbf {K} \otimes_ {\mathbf {A}} \mathrm{H} _ {p} (\mathbf {C})) \rightarrow \mathrm{H} _ {p} (\mathbf {K} \otimes_ {\mathbf {A}} \mathbf {C}) \rightarrow \mathrm{H} _ {1} (\mathbf {K} \otimes_ {\mathbf {A}} \mathrm{H} _ {p - 1} (\mathbf {C})) \rightarrow 0.
\]

\begin{enumerate}
\def\labelenumi{\alph{enumi})}
\setcounter{enumi}{1}
\tightlist
\item
  若 K 是投射的，则对每个 \(p \in \mathbf{N}\) ，我们有正合列
\end{enumerate}

\[
\begin{array}{r l} 0 \to \mathrm{H} ^ {1} (\operatorname{Homgr} _ {\mathbf {A}} (\mathbf {K}, \mathrm{H} ^ {p - 1} (\mathbf {C}))) & \to \mathrm{H} ^ {p} (\operatorname{Homgr} _ {\mathbf {A}} (\mathbf {K}, \mathbf {C})) \\ & \to \mathrm{H} ^ {0} (\operatorname{Homgr} _ {\mathbf {A}} (\mathbf {K}, \mathrm{H} ^ {p} (\mathbf {C}))) \to 0. \end{array}
\]

我们来证明 \(a\) ， \(b\) ) 的证明类似。对每个 \(i\) ，记 \(\mathbf{K}_{(i)}\) 为复形 \(\mathbf{K}_i(-i)\) 。我们有一个复形的正合列，它作为 A-模的序列是分裂的

\[
0 \rightarrow \mathrm{K} _ {(0)} \xrightarrow {\alpha} \mathrm{K} \xrightarrow {\beta} \mathrm{K} _ {(1)} \rightarrow 0;
\]

序列

\[
0 \longrightarrow \mathrm{K} _ {(0)} \otimes_ {\mathrm{A}} \mathrm{C} \xrightarrow {\alpha \otimes 1} \mathrm{K} \otimes_ {\mathrm{A}} \mathrm{C} \xrightarrow {\beta \otimes 1} \mathrm{K} _ {(1)} \otimes_ {\mathrm{A}} \mathrm{C} \longrightarrow 0\tag{14}
\]

是正合的，并且由于 K 是平坦的，同态

\[
\gamma_ {0, p} (\mathbf {K} _ {(0)}, \mathbf {C}): \mathbf {K} _ {(0)} \otimes_ {\mathbf {A}} \mathrm{H} _ {p} (\mathbf {C}) \rightarrow \mathrm{H} _ {p} (\mathbf {K} _ {(0)} \otimes_ {\mathbf {A}} \mathbf {C})
\]

\[
\gamma_ {1, p - 1} (\mathbf {K} _ {(1)}, \mathbf {C}): \mathbf {K} _ {(1)} \otimes_ {\mathbf {A}} \mathrm{H} _ {p - 1} (\mathbf {C}) \rightarrow \mathrm{H} _ {p} (\mathbf {K} _ {(1)} \otimes_ {\mathbf {A}} \mathbf {C})
\]

是双射的（X，第 66 页，推论 2）。让我们计算连接同态 \(\partial (\alpha \otimes 1, \beta \otimes 1)\) ；根据定义，它把闭链 \(\sum a_i \otimes b_i\) 的类映到 \(\sum da_i \otimes b_i\) 的类，这意味着

\[
\partial (\alpha \otimes 1, \beta \otimes 1) \circ \gamma (\mathbf {K} _ {(1)}, \mathbf {C}) = \gamma (\mathbf {K} _ {(0)}, \mathbf {C}) \circ (d _ {\mathbf {K}} \otimes 1).
\]

因此，与 (14) 相关联的同调正合列具有如下形式：

\[
\begin{array}{r l} \mathrm{K} _ {1} \otimes \mathrm{H} _ {p} (\mathrm{C}) & \xrightarrow {d _ {\mathbf {k}} \otimes 1} \mathrm{K} _ {0} \otimes \mathrm{H} _ {p} (\mathrm{C}) \longrightarrow \mathrm{H} _ {p} (\mathrm{K} \otimes \mathrm{C}) \\ & \longrightarrow \mathrm{K} _ {1} \otimes \mathrm{H} _ {p - 1} (\mathrm{C}) \xrightarrow {d _ {\mathbf {k}} \otimes 1} \mathrm{K} _ {0} \otimes \mathrm{H} _ {p - 1} (\mathrm{C}), \end{array}
\]

由此得 a）。

将引理 3 的 a) 部分应用于复形 \(\mathbf{K}(u)\) ，并利用交换同构，我们得到：

命题 4. --- 对每个复形 C，我们有正合列。

\[
0 \rightarrow \mathrm{A} / x \mathrm{A} \otimes_ {\mathrm{A}} \mathrm{H} _ {p} (\mathrm{C}) \rightarrow \mathrm{H} _ {p} (x, \mathrm{C}) \rightarrow \operatorname{Hom} _ {\mathrm{A}} (\mathrm{A} / x \mathrm{A}, \mathrm{H} _ {p - 1} (\mathrm{C})) \rightarrow 0.
\]

推论 1. --- \(\mathbf{H}_{p}(x,\mathbf{C})=0\) 当且仅当比率为 \(x\) 的位似到 \(\mathbf{H}_{p}(\mathbf{C})\) 上是满射的，且比率为 \(x\) 的位似到 \(\mathbf{H}_{p-1}(\mathbf{C})\) 上是单射的。

推论 2. --- 设 \(x = (x_1, \ldots, x_n)\) 是 A 中元素的一个序列，M 是一个 A-模， \(x'\) 为序列 \((x_1, \ldots, x_{n-1})\) 。我们有正合列

\[
0 \rightarrow \mathrm{A} / x _ {n} \mathrm{A} \otimes_ {\mathrm{A}} \mathrm{H} _ {p} (x ^ {\prime}, \mathrm{M}) \rightarrow \mathrm{H} _ {p} (x, \mathrm{M}) \rightarrow \operatorname{Hom} _ {\mathrm{A}} (\mathrm{A} / x _ {n} \mathrm{A}, \mathrm{H} _ {p - 1} (x ^ {\prime}, \mathrm{M})) \rightarrow 0.
\]

推论 3. --- \(\mathrm{H}_{i}(\mathbf{x},\mathbf{M})\) 对所有 i \textgreater{} 0 均为零，当且仅当比率为 \(x_{n}\) 的位似到 \(\mathrm{H}_{i}(\mathbf{x}^{\prime},\mathbf{M})\) 上对 i \textgreater{} 0 是双射，并且比率为 \(x_{n}\) 的位似到 \(\mathrm{M}/(\mathbf{x}^{\prime})\) M 上是单射。

\subsection{完全切割族}

设 A 是一个环，M 是一个 A-模， \(x = (x_{i})_{i \in I}\) 是 A 中元素的一个族。

定义 2. --- 称族 x 对 M 是完全切割的，如果有 \(\mathrm{H}_{i}(\mathbf{x},\mathbf{M})=0\) 对 i\textgreater0 。

若 I 是有限的，则这等价于（X，第 150 页）说有 \(\mathbf{H}^i (\pmb {x},\mathbf{M}) = 0\) 对 \(i <   \mathrm{Card}\) (I)。

下面的命题使我们能够给出完全切割族的一些例子。

命题 5. --- 设 \(x = (x_1, \ldots, x_n)\) 是 A 中元素的一个序列。如果对 \(i=1, \ldots, n\) ，比率为 \(x_i\) 的位似在 A-模 \(M/(x_1 M + \cdots + x_{i-1} M)\) 中是单射的，则序列 \(x\) 对 M 是完全切割的。

满足命题条件的序列称为对 M 正则的，或 M-正则的。注意，此性质一般依赖于 \(x_i\) 的顺序。例如，序列 (1, 0) 总是 M-正则的，而序列 (0, 1) 仅当 M 为零模时才是 M-正则的。另一方面，一个序列是否完全切割并不依赖于各项的顺序。

让我们对 \(n\) 作归纳来证明该命题， \(n=0\) 的情形是显然的。设 \(x'=(x_{1},...,x_{n-1})\) ；若序列 \(x\) 是 M-正则的，则序列 \(x'\) 是 M-正则的，且以 \(x_{n}\) 乘 \(\mathrm{M}/(\mathbf{x}')\mathrm{M}\) 的位似是单射；由归纳假设，我们有 \(\mathrm{H}_{i}(\mathbf{x}',\mathrm{M})=0\) 对 \(i>0\) ；于是由 X，第 157 页，推论 3，可知 \(\mathrm{H}_{i}(\mathbf{x},\mathrm{M})=0\) 对 \(i>0\) 。

例. --- 设 \(k\) 是一个环；取 \(\mathbf{A} = k[\mathbf{X}_1, \ldots, \mathbf{X}_n]\) 和 \(\boldsymbol{x} = (\mathbf{X}_1, \ldots, \mathbf{X}_n)\) 。序列 \(\boldsymbol{x}\) 是 \(\mathbf{A}\) -正则的，且该命题重新给出次数 \(>0\) 时复形 \(\mathbf{S}_k(k^n) \otimes_k \boldsymbol{\Lambda}_k(k^n)\) 的零调性（参见 X，第 152 页，命题 3）。

同样地，在形式幂级数环 \(\hat{\mathbf{A}} = k[[\mathbf{X}_1, \ldots, \mathbf{X}_n]]\) 中，序列 \((\mathbf{X}_1, \ldots, \mathbf{X}_n)\) 是 \(\hat{\mathbf{A}}\) -正则的，因此对 \(\hat{\mathbf{A}}\) 完全切割。

命题 6. --- a) 若 \(\sum_{i\in I}x_i\mathbf{A} = \mathbf{A}\) ，则族 \((x_i)_{i\in I}\) 对 M 是完全切割的。

\begin{enumerate}
\def\labelenumi{\alph{enumi})}
\setcounter{enumi}{1}
\tightlist
\item
  设 \(x = (x_1, \ldots, x_n)\) 是 A 的元素序列。设 \((a_{ij}) \in \mathbf{GL}_n(\mathbf{A})\) ；令
\end{enumerate}

\[
y _ {i} = \sum_ {j} a _ {i j} x _ {j}.
\]

若序列 x 对 M 完全切割，则序列 \((y_{1}, \ldots, y_{n})\) 亦然。c) 设 \(k_{1}, \ldots, k_{n}\) 为整数 \(\geqslant 1\) ；序列 \((x_{1}^{k_{1}}, \ldots, x_{n}^{k_{n}})\) 对 M 完全切割，当且仅当序列 x 对 M 完全切割。断言 a) 由第 148 页推论 1 得出。

我们来证明 \(b\) )。设 \(f: \mathbf{A}^n \to \mathbf{A}^n\) 为由矩阵 \(^t(a_{ij})\) 定义的同构；由 X，第 151 页可知 \(1_{\mathbf{M}} \otimes \Lambda(f)\) 是复形 \(\mathbf{K}.(y, \mathbf{M})\) 到复形 \(\mathbf{K}.(x, \mathbf{M})\) 上的同构，从而得 \(b\) 。

为了证明 \(c\) ），显然只需证明：当 \(k\) 是整数 \(\geqslant 1\) 时，序列 \((x_{1}, \ldots, x_{n-1}, x_{n}^{k})\) 对 M 完全切割当且仅当序列 \(x\) 对 M 完全切割。设 \(x' = (x_{1}, \ldots, x_{n-1})\) ；由第 157 页推论 3，第一个条件（相应地，第二个条件）意味着比率为 \(x_{n}^{k}\) （相应地， \(x_{n}\) ）的位似在 H \(_{i}(x', M)\) 中对 \(i \geqslant 1\) 是双射的，且在 M/(x') M 中是单射的。这两个条件显然等价，因此得 \(c\) ）。

注记。--- 1) 对正则序列，与 c) 类似的断言成立（X，第 207 页，习题 5）。

命题 7. --- a) 设 N 是一个平坦的 A-模。如果族 x 对 M 完全切割，则它对 M ⊗\_A N 也完全切割。

\begin{enumerate}
\def\labelenumi{\alph{enumi})}
\setcounter{enumi}{1}
\tightlist
\item
  设 \(0 \to \mathbf{M}' \to \mathbf{M} \to \mathbf{M}'' \to 0\) 是 A-模的一个正合列。如果族 \(x\) 对 \(\mathbf{M}'\) 和 \(\mathbf{M}''\) 都完全切割，则它对 \(\mathbf{M}\) 也完全切割。
\end{enumerate}

复形 \(\mathbf{K}.\left(\mathbf{x},\mathbf{M}\otimes_{\mathbf{A}}\mathbf{N}\right)\) 由定义同构于 \(\mathbf{K}.\left(\mathbf{x},\mathbf{M}\right)\otimes_{\mathbf{A}}\mathbf{N}\) ；由于 N 是平坦的，我们推出一个同构 H. \((\mathbf{x},\mathbf{M})\otimes_{\mathbf{A}}\mathbf{N}\to\mathrm{H}.\left(\mathbf{x},\mathbf{M}\otimes_{\mathbf{A}}\mathbf{N}\right)\) （X，第 66 页，推论 2），由此得 \(a\) ）。

由于复形 \(\mathbf{K}.\left(\mathbf{x}\right)\) 是平坦的，我们有一个复形的正合列

\[
0 \rightarrow \mathbf {K}. (x, \mathrm{M} ^ {\prime}) \rightarrow \mathbf {K}. (x, \mathrm{M}) \rightarrow \mathbf {K}. (x, \mathrm{M} ^ {\prime \prime}) \rightarrow 0;
\]

断言 b) 由相伴的同调正合列得出。

注记。--- 2) 对正则序列而言，与 a) 和 b) 类似的断言是直接的。

\begin{enumerate}
\def\labelenumi{\arabic{enumi})}
\setcounter{enumi}{2}
\tightlist
\item
  如果族 \(x\) 对 \(A\) 完全切割，则复形 \(K.(x, A)\) 定义了 \(A\) -模 \(A/x\) 的一个自由分解，其中 \(x = \sum_{i \in I} x_i A\) ；因此对每个整数 \(i \geqslant 0\) 以及每个 \(A\) -模 \(M\) 我们有同构
\end{enumerate}

\[
\operatorname{Ext} _ {\mathbf {A}} ^ {n - i} (\mathbf {A} / \mathfrak {x}, \mathbf {M}) \rightarrow \mathrm{H} ^ {n - i} (\mathfrak {x}, \mathbf {M}) \rightarrow \mathrm{H} _ {i} (\mathfrak {x}, \mathbf {M}) \rightarrow \operatorname{Tor} _ {i} ^ {\mathbf {A}} (\mathbf {A} / \mathfrak {x}, \mathbf {M}).\tag{15}
\]

\begin{enumerate}
\def\labelenumi{\arabic{enumi})}
\setcounter{enumi}{3}
\tightlist
\item
  我们称序列 \(\boldsymbol{x} = (x_1, \dots, x_n)\) 是 M-余正则的，如果（记 \((x_i)_M\) 为 M 中比率为 \(x_i\) 的位似）模中比率为 \(x_i\) 的位似
\end{enumerate}

\[
\operatorname{Ker} \left(x _ {1}\right) _ {\mathbf {M}} \cap \dots \cap \operatorname{Ker} \left(x _ {i - 1}\right) _ {\mathbf {M}}
\]

对 \(i = 1, \ldots, n\) 是满射的。于是我们有 \(\mathbf{H}_{i}(\mathbf{x}, \mathbf{M}) = 0\) 对 \(i < n\) ；其证明方式与命题 5 相同。

\[
\mathrm{A} = k \left[ \mathrm{D} _ {1}, \dots , \mathrm{D} _ {n} \right],
\]

\[
\mathbf {Q} \text {-algebra},
\]

\(M = k[T_1, \ldots, T_n]\) ，赋予其一个 A-模结构，使得 \(D_i P = \partial P / \partial T_i\) 对所有 \(P \in M (1 \leqslant i \leqslant n)\) 成立。立即可验证序列 \((D_1, \ldots, D_n)\) 是 M-余正则的；结合第 155 页的例子，可以推出 \(k[T_1, \ldots, T_n]\) 在 \(k\) 上的德拉姆复形在次数 \textgreater{} 0 时是零调的。

\subsection{完全切割序列的一个判据}

设 A 为一个环，M 为一个 A-模，x 为 A 的一个理想。M 上的 x-进拓扑是与 M 的群结构相容的拓扑，且在该拓扑下，子模 \(x^{r}\) M ( \(r \geqslant 0\) ) 的集合构成零的一个基本邻域系（TG，III，第 5 页，例）。该拓扑是分离的当且仅当

\[
\bigcap_ {r \geqslant 0} x ^ {r} = 0.
\]

现在假设理想 x 由 A 中元素的序列 \(\boldsymbol{x}=(x_{1},\ldots,x_{n})\)

生成。考虑分次 A-模 \(\bigoplus_{r\geqslant 0}x^{r}\) M 以及 0 次的分次 A-同态

\[
a _ {\mathbf {M}} ^ {\mathbf {x}}: \mathrm{A} [ \mathrm{X} _ {1}, \dots , \mathrm{X} _ {n} ] \otimes_ {\mathbf {A}} \mathrm{M} \rightarrow \bigoplus_ {r \geqslant 0} \mathfrak {x} ^ {r} \mathrm{M}
\]

使得 \(a_{\mathbf{M}}^{\star}(\mathbf{P} \otimes m) = \mathbf{P}(x_1, \ldots, x_n) \, m\) ，如果 \(\mathbf{P}\) 是 \(\mathbf{X}_1, \ldots, \mathbf{X}_n\) 中的齐次多项式且 \(m \in \mathbf{M}\) 。设 \(\mathfrak{d}\) 为 \(\mathbf{A}[\mathbf{X}_1, \ldots, \mathbf{X}_n]\) 中由元素 \((x_i \, \mathbf{X}_j - x_j \, \mathbf{X}_i)\) （ \(1 \leqslant i < j \leqslant n\) ）生成的理想。我们有 \(\mathbf{P}(x_1, \ldots, x_n) = 0\) ，如果 \(\mathbf{P} \in \mathfrak{d}\) ，从而 \(a_{\mathbf{M}}^{\star}\) 通过过渡到商给出一个 0 次的分次 A-同态

\[
\alpha_ {\mathrm{M}} ^ {x}: \left(\mathrm{A} \left[ \mathrm{X} _ {1}, \dots , \mathrm{X} _ {n} \right] / \mathfrak {d}\right) \otimes_ {\mathrm{A}} \mathrm{M} \rightarrow \bigoplus_ {r \geqslant 0} x ^ {r} \mathrm{M}.
\]

通过与 A/x 作张量积，我们从 \(\alpha_{\mathrm{M}}^{\mathbf{x}}\) 导出一个 0 次的分次 A-同态

\[
\beta_ {\mathrm{M}} ^ {x}: (\mathrm{A} / \mathfrak {x}) [ \mathrm{X} _ {1}, \dots , \mathrm{X} _ {n} ] \otimes_ {\mathrm{A}} \mathrm{M} \rightarrow \bigoplus_ {r \geqslant 0} (\mathfrak {x} ^ {r} \mathrm{M} / \mathfrak {x} ^ {r + 1} \mathrm{M}).
\]

同态 \(a_{\mathbf{M}}^{x}\) 、\(\alpha_{\mathbf{M}}^{x}\) 与 \(\beta_{\mathbf{M}}^{x}\) 是满射的。

定理 1. --- 考虑以下条件：

\begin{enumerate}
\def\labelenumi{(\roman{enumi})}
\item
  序列 x 是 M-正则的（X，第 158 页）。
\item
  序列 x 对 M 完全切割（X，第 157 页，定义 2）。
\item
  我们有 \(\mathbf{H}_1(x, \mathbf{M}) = 0\) 。
\item
  \(L^{\prime}\) 同态 \(\alpha_{\mathbf{M}}^{x}:(\mathbf{A}[\mathbf{X}_{1},\dots,\mathbf{X}_{n}]/\mathfrak{d})\otimes_{\mathbf{A}}\mathbf{M}\to\bigoplus x^{r}\mathbf{M}\) 是双射。
\item
  同态 \(\beta_{\mathbf{M}}^{\mathbf{x}}:(\mathbf{A}/\mathfrak{x})[\mathbf{X}_{1},..., \mathbf{X}_{n}]\otimes_{\mathbf{A}}\mathbf{M}\to\bigoplus_{r\geqslant0}(\mathfrak{x}^{r}\mathbf{M}/\mathfrak{x}^{r+1}\mathbf{M})\) 是双射。
\end{enumerate}

那么我们有蕴含关系 (i)⇒(ii)⇒(iii)⇔(iv)⇒(v)。若对 \(1 \leqslant i \leqslant n\) ，A-模 \(\mathrm{M}/(x_{1}\mathrm{M} + \cdots + x_{i-1}\mathrm{M})\) 对 x-进拓扑是分离的，则条件 (i) 至 (v) 等价。

该定理将在第 8 节和第 9 节中证明。

推论 1. --- 若 A 是诺特环，若 A-模 M 是有限生成的，且若 \(x_{i}\) 属于 A 的根，则定理的条件 (i) 至 (v) 等价。事实上，在每个模 \(\mathbf{M}/(x_{1} \mathbf{M} + \cdots + x_{i-1} \mathbf{M})\) 上 x-进拓扑都是分离的（AC III，§ 3，第 3 号，命题 6）。

推论 2. --- 设 A 是一个正分次环，M 是一个下有界的分次 A-模，且 \(x_i\) 为 A 中次数 \textgreater{} 0 的齐次元素。则定理的条件 (i) 至 (v) 等价。

对任何下有界的分次 A-模 N，x-进拓扑确实是分离的，因为若 \(N_{n}=0\) 对 \(n<n_{0}\) ，则我们有 \(x^{a}N\subset\sum_{j\geqslant n_{0}+a}N_{j}\) 对所有 \(a\geqslant0\) 。

推论 3. --- 假设模 \(\mathbf{M}/(x_1\mathbf{M} + \cdots + x_{i-1}\mathbf{M})\) 对 \(\mathbf{x}\) -进拓扑是分离的（ \(1 \leqslant i \leqslant n\) ）；设 \(p\) 为介于 1 与 \(n\) 之间的整数。序列 \(\mathbf{x}\) 对 \(\mathbf{M}\) 完全切割，当且仅当序列 \((x_1, \ldots, x_p)\) 对 \(\mathbf{M}\) 完全切割且序列 \((x_{p+1}, \ldots, x_n)\) 对 \(\mathbf{M}/(x_1\mathbf{M} + \cdots + x_p\mathbf{M})\) 完全切割。

事实上，如果在陈述中将“完全切割序列”替换为“正则序列”，则该推论是显然的；而根据定理，这两个概念在此处是一致的。

注记。--- 设 \(x = (x_1, \ldots, x_n)\) 是 A 中元素的序列，使得 \(H_1(x, A) = 0\) ；则满射同态 \(u: A^n \to x\) （它使得 \(u(\sum a_i e_i) = \sum a_i x_i\) ）的核由元素 \(X_j e_i - X_i e_j\) 生成；因此，A-代数 \(A[X_1, \ldots, X_n] / \mathfrak{d}\) 同构于对称代数 \(S_A(x)\) （III，第 69 页，命题 4）。于是由定理 1 可知，代数同态 \(S_A(x) \to \bigoplus_n x^n\) （它由 \(x\) 到 \(\bigoplus_n x^n\) 中的典范单射导出）是一个同构。对同态 \(S_A(x / x^2) \to \bigoplus_n x^n / x^{n+1}\) 同样成立。

\subsection{定理1的证明：第一部分}

蕴含关系 (i) \(\Rightarrow\) (ii) 已在（X，第 157 页，命题 5）中证明。蕴含关系 (ii) \(\Rightarrow\) (iii) 是显然的；对 (iv) \(\Rightarrow\) (v) 同理，因为 \(\beta_{\mathbf{M}}^{x}\) 与 \(\alpha_{\mathbf{M}}^{x} \otimes 1_{\mathbf{A}/\mathfrak{x}}\) 等同。

为证明 (iv) 蕴含 (iii)，考虑在 1 次分量上诱导的同态 \((\alpha_{M}^{x})_{1}\) 。设 E 表示分次 A-模 \(A[X_{1},\ldots,X_{n}]\) ；A-模 \(E_{1}\) 以 \(X_{1},\ldots,X_{n}\) 为基，而 \(\mathfrak{d}_{1}\) 是 \(E_{1}\) 的子 A-模，由元素 \((x_{i}X_{j}-x_{j}X_{i})\) （ \(1\leqslant i<j\leqslant n\) ）生成。因此 \(((E/\mathfrak{d})\otimes_{A}M)_{1}\) 与 \(K_{1}(x,M)/B_{1}(K.(x,M))\) 等同，同态 \((\alpha_{M}^{x})_{1}\) 等同于从 \(K_{1}(x,M)/B_{1}(K.(x,M))\) 到 \(B_{0}(K.(x,M))\) 上的、由 \(d_{1}\) 诱导的映射。因此 \(H_{1}(x,M)\) 的消失等价于 \((\alpha_{M}^{x})_{1}\) 的单射性，从而得到蕴含关系 (iv) \(\Rightarrow\) (iii)。

为证明(iii)蕴含(iv)，我们将使用以下引理：

引理 4. --- 设 \(\mathbf{A}_{0}\) 为环 \(\mathbf{Z}[\mathrm{T}_{1},\ldots,\mathrm{T}_{n}]\) ，并设 \(u:\mathbf{A}_{0}[\mathrm{X}_{1},\ldots,\mathrm{X}_{n}] \to \mathbf{A}_{0}[\mathrm{U}]\) 为 \(\mathbf{A}_{0}\) -代数的、使得 \(u(\mathbf{X}_{i}) = \mathbf{T}_{i}\) U 的同态。\(u\) 的核是理想 \(\mathfrak{d}_{0}\) ，它由 \(\mathbf{A}_{0}[\mathbf{X}_{1},\ldots,\mathbf{X}_{n}]\) 中形如（ \(\mathrm{T}_{i}\mathbf{X}_{j}-\mathrm{T}_{j}\mathbf{X}_{i}\) ）的元素（ \(1 \leqslant i < j \leqslant n\) ）生成。若 t 是 \(\mathbf{A}_{0}\) 中由 ( \(\mathrm{T}_{1},\ldots,\mathrm{T}_{n}\) ) 生成的理想，则 \(u\) 诱导出同构

\[
\overline {{u}}: \mathrm{A} _ {0} [ \mathrm{X} _ {1}, \dots , \mathrm{X} _ {n} ] / \mathfrak {d} _ {0} \rightarrow \bigoplus_ {r \geqslant 0} t ^ {r}.
\]

显然只需证明第一个断言。对每个自然整数序列 \(\alpha = (\alpha_{1},\ldots,\alpha_{n})\) 及每个整数 \(k\in[0,n]\) ，记 \(\mathbf{P}_{\alpha,k}\) 为单项式

\[
\mathrm{T} _ {1} ^ {\alpha_ {1}} \dots \mathrm{T} _ {k} ^ {\alpha_ {k}} \mathrm{X} _ {k + 1} ^ {\alpha_ {k + 1}} \dots \mathrm{X} _ {n} ^ {\alpha_ {n}};
\]

令 N 为 \(A_{0}[X_{1}, \ldots, X_{n}]\) 的子 Z-模，由 \(P_{\alpha,k}\) 对 \(\alpha \in N^{n}\) 和 \(0 \leqslant k \leqslant n\) 生成。我们将证明 u 在 N 上的限制是单射，且 \(A_{0}[X_{1}, \ldots, X_{n}] = \mathfrak{d}_{0} + N\) ；由于有 \(d_{0} \subset Ker u\) ，由此即得引理。

注意 N 由集合 S 生成，S 由 \(P_{0,0} = 1\) 以及那些 \(P_{\alpha,k}\) （对它们有 \(\alpha_k \neq 0\) ）组成。为证明 u 在 N 上的限制是单射的，只需证明 S 中两个不同的元素在 u 下的像是 \(A_0[U]\) 中不同的单项式。现在有 \(u(P_{\alpha,k}) = T^\alpha U^{\sum_{i \geq k} \alpha_i}\) ，因此等式 \(u(P_{\alpha,k}) = u(P_{\alpha',k'})\) 蕴含 \(\alpha = \alpha'\) 与 \(\sum_{i \geq k} \alpha_i = \sum_{i \geq k'} \alpha_i\) 。假设 \(P_{\alpha,k}\) 和 \(P_{\alpha',k'}\) 属于 S。若 \(\alpha = 0\) ，则有 \(k = k' = 0\) ；若 \(\alpha \neq 0\) ，我们得到 \(k = k'\) ，因为 \(\alpha_k\) 与 \(\alpha_{k'}\) 非零，从而结论成立。

让我们证明每个单项式 \(\mathbf{T}^{\alpha}\mathbf{X}^{\beta}\in\mathbf{A}_{0}[\mathbf{X}_{1},..., \mathbf{X}_{n}]\) 模 \(\mathfrak{d}_{0}\) 同余于某个 \(\mathbf{P}_{\eta,k}\) 。对每个序列 \(\lambda\in\mathbf{N}^{n}\) ，我们用 \(i(\lambda)\) （相应地， \(j(\lambda)\) ）表示最小的（相应地，最大的）整数 \(k\in[1,n]\) ，它使得 \(\lambda_{k}\neq0\) 。在单项式 \(\mathbf{T}^{\gamma}\mathbf{X}^{\delta}\) （它们与 \(\mathbf{T}^{\alpha}\mathbf{X}^{\beta}\) 模 \(\mathfrak{d}_{0}\) 同余）的集合中，我们选择一个使有理整数 \(j(\gamma)-i(\delta)\) 为最小的；让我们证明此时有 \(j(\gamma)-i(\delta)<0\) 。假设我们有 \(j(\gamma)\geqslant i(\delta)\) ；记 \(j=j(\gamma)\) , \(i=i(\delta)\) ，并设 \(\varepsilon=\inf(\gamma_{j},\delta_{i})\) 。由于 \((\mathbf{T}_{j}^{\varepsilon}\mathbf{X}_{i}^{\varepsilon}-\mathbf{T}_{i}^{\varepsilon}\mathbf{X}_{j}^{\varepsilon})\) 能被 \((\mathbf{T}_{j}\mathbf{X}_{i}-\mathbf{T}_{i}\mathbf{X}_{j})\) 整除，故属于 \(\mathfrak{d}_{0}\) ，我们看到 \(\mathbf{T}^{\gamma}\mathbf{X}^{\delta}\) 模 \(\mathfrak{d}_{0}\) 同余于 \(\mathbf{T}^{\gamma'}\mathbf{X}^{\delta'}\) ，其中 \(\gamma_{i}'=\gamma_{i}+\varepsilon\) , \(\gamma_{j}'=\gamma_{j}-\varepsilon\) , \(\gamma_{k}'=\gamma_{k}\) （ \(k\neq i,j\) ），且 \(\delta_{i}'=\delta_{i}-\varepsilon\) , \(\delta_{j}'=\delta_{j}+\varepsilon\) , \(\delta_{k}'=\delta_{k}\) （ \(k\neq i,j\) ）。由于 \(\gamma_{j}'\) 或 \(\delta_{i}'\) 为零，我们有 \(j(\gamma')-i(\delta')<j(\gamma)-i(\delta)\) ，这与 \(j(\gamma)-i(\delta)\) 的极小性相矛盾。

因此我们有 \(j(\gamma) < i(\delta)\) ，从而 \(\mathbf{T}^{\gamma}\mathbf{X}^{\delta} \in \mathbf{N}\) ，这就完成了引理的证明。

让我们证明 (iii) 蕴含 (iv)。考虑环 \(\mathbf{A}_0 = \mathbf{Z}[\mathrm{T}_1, \ldots, \mathrm{T}_n]\) 以及 \(\mathbf{A}_0\) 中由 \(\mathrm{T}_1, \ldots, \mathrm{T}_n\) 生成的理想 t。赋予 M 一个 \(\mathbf{A}_0\) -模结构，使得 \(\mathrm{T}_i m = x_i m\) 对 \(m \in \mathbf{M}\) , \(1 \leqslant i \leqslant n\) 成立。根据 X，第 155 页， \(\mathrm{H}_i(x, \mathbf{M})\) 与 \(\mathbf{H}_i(\mathbf{T}, \mathbf{M})\) 典范地等同。

由于序列 \(\mathbf{T}\) 对 \(\mathbf{A}_{0}\) 是正则的，由蕴含关系 (i) \(\Rightarrow\) (ii) 可知，复形 \(\mathbf{K}.\left(\mathbf{T},\mathbf{A}_{0}\right)\) 定义了 \(\mathbf{A}_{0}\) -模 \(\mathbf{A}_{0}/t\) 的一个自由分解。注意，该模与 \(\mathbf{Z}\) 等同，配备 \(\mathbf{A}_{0}\) -模结构，使得 \(\mathrm{T}_{i}\mathbf{Z}=0\) 对 \(1\leqslant i\leqslant n\) 成立。因此条件 (iii) 等价于 \(\operatorname{Tor}_{1}^{\mathbf{A}_{0}}(\mathbf{M},\mathbf{Z})=0\) 。

让我们证明 (iii) 蕴含 \(\mathrm{Tor}_{1}^{\mathbf{A}_{0}}(\mathbf{M},\mathbf{A}_{0}/t^{r})=0\) 对所有 \(r\geqslant1\) 成立。这对 \(r=1\) 的情形由前面的论述得出。在一般情形，考虑正合列

\[
0 \rightarrow t ^ {r} / t ^ {r + 1} \rightarrow A _ {0} / t ^ {r + 1} \rightarrow A _ {0} / t ^ {r} \rightarrow 0.\tag{16}
\]

该 \(A_{0}\) -模 \(t^{r}/t^{r+1}\) 同构于 Z 的有限份拷贝的直积；因此我们有 \(\operatorname{Tor}_{1}^{\mathbf{A}_{0}}(\mathbf{M}, t^{r}/t^{r+1}) = 0\) 。由此，通过对 r 归纳，我们推出

\[
\operatorname{Tor} _ {1} ^ {\mathrm{A} _ {0}} \left(\mathrm{M}, \mathrm{A} _ {0} / t ^ {r}\right) = 0 \quad \text { for   all } r.
\]

因此，正合列 (16) 通过与 M 作张量积，给出一个正合列

\[
0 \rightarrow (t ^ {r} / t ^ {r + 1}) \otimes_ {\mathrm{A} _ {0}} \mathrm{M} \rightarrow \mathrm{M} / x ^ {r + 1} \mathrm{M} \rightarrow \mathrm{M} / x ^ {r} \mathrm{M} \rightarrow 0\tag{17}
\]

从而得到 \((\mathfrak{t}^r /\mathfrak{t}^{r + 1})\otimes_{\mathbf{A}_0}\mathbf{M}\) 到 \(\mathfrak{x}^r\mathbf{M} / \mathfrak{x}^{r + 1}\mathbf{M}\) 上的一个同构。考虑正合列 \(0\to \mathfrak{t}^{r + 1}\to \mathfrak{t}^r\to \mathfrak{t}^r /\mathfrak{t}^{r + 1}\to 0\) ，然后对 \(r\) 归纳可以看出，映射 \(m_r:\mathfrak{t}^r\otimes_{\mathbf{A}_0}\mathbf{M}\rightarrow \mathfrak{x}^r\mathbf{M}\) （它由 \(\mathbf{A}_0\) 在 \(\mathbf{M}\) 上的作用诱导）是一个同构。

为证明 (iv)，只需观察到图表

\[
\begin{array}{c} \left(\mathrm{A} _ {0} \left[ \mathrm{X} _ {1}, \dots , \mathrm{X} _ {n} \right] / \mathfrak {d} _ {0}\right) \otimes_ {\mathrm{A} _ {0}} \mathrm{M} \xrightarrow {\bar {u} \otimes 1 _ {\mathrm{M}}} \bigoplus_ {r \geqslant 0} \left(\mathrm{t} ^ {r} \otimes_ {\mathrm{A} _ {0}} \mathrm{M}\right) \\ \Bigg \downarrow^ {e} \\ ^ {\prime} (\mathrm{A} [ \mathrm{X} _ {1}, \dots , \mathrm{X} _ {n} ] / \mathfrak {d}) \otimes_ {\mathrm{A}} \mathrm{M} \xrightarrow {\alpha_ {\mathrm{M}} ^ {*}} \bigoplus_ {r \geqslant 0} x ^ {r} \mathrm{M} \end{array}
\]

其中 e 是纯量扩张的典范同构（II，第 85 页，命题 6），该图表是交换的，然后应用引理 4 即可。

\subsection{定理1的证明：第二部分}

再次考虑正合列

\[
0 \longrightarrow t ^ {r} / t ^ {r + 1} \xrightarrow {i _ {r}} A _ {0} / t ^ {r + 1} \xrightarrow {p _ {r}} A _ {0} / t ^ {r} \longrightarrow 0\tag{16}
\]

以及相伴的挠模正合列

\[
\begin{array}{l} \text {Tor} _ {1} ^ {\mathrm{A} _ {0}} (\mathrm{A} _ {0} / t ^ {r + 1}, \mathrm{M}) \longrightarrow \mathrm{Tor} _ {1} ^ {\mathrm{A} _ {0}} (\mathrm{A} _ {0} / t ^ {r}, \mathrm{M}) \\ \longrightarrow (t ^ {r} / t ^ {r + 1}) \otimes_ {\mathrm{A} _ {0}} \mathrm{M} \xrightarrow {i , \otimes 1 _ {\mathrm{M}}} (\mathrm{A} _ {0} / t ^ {r + 1}) \otimes_ {\mathrm{A} _ {0}} \mathrm{M} \xrightarrow {p , \otimes 1 _ {\mathrm{M}}} (\mathrm{A} _ {0} / t ^ {r}) \otimes_ {\mathrm{A} _ {0}} \mathrm{M} \longrightarrow 0. \end{array} \tag {18}
\]

\(p_r \otimes 1_M\) 的核与 \(x^r M/x^{r+1} M\) 等同；此外， \(t^r/t^{r+1}\) 被 \(T_i\) 零化，并与次数为 \(r\) 的 \(A_0\) 的齐次分量等同，因此同态 \((t^r/t^{r+1}) \otimes_{A_0} M \to x^r M/x^{r+1} M\) （它由 \(i_r \otimes 1_M\) 导出）在次数 \(r\) 处与同态 \(\beta_M^x\) 的齐次分量等同。因此由正合列（18）可知，条件（v）等价于

（v'）：同态 \(\operatorname{Tor}_{1}^{\mathbf{A}_{0}}(p_{r},1_{\mathbf{M}})\) : \(\operatorname{Tor}_{1}^{\mathbf{A}_{0}}(\mathbf{A}_{0}/t^{r+1},\mathbf{M})\to\operatorname{Tor}_{1}^{\mathbf{A}_{0}}(\mathbf{A}_{0}/t^{r},\mathbf{M})\) 对所有 \(r\geqslant1\) 是满射的。

我们仍需证明蕴含关系 (v) \(\Rightarrow\) (i) 在诸模

\[
\mathbf {M} / (x _ {1} \mathbf {M} + \dots + x _ {i - 1} \mathbf {M})
\]

对 \(\mathfrak{x}\) -进拓扑是分离的时成立（ \(1 \leqslant i \leqslant n\) ）。设 \(\overline{\mathbf{M}}\) 为 A-模 \(\mathbf{M}/x_1 \mathbf{M}\) 。根据定义，序列 \(\mathbf{x}\) 对 \(\mathbf{M}\) 是正则的当且仅当 \((x_1)_\mathbf{M}\) 是单射且序列 \(\mathbf{x}' = (x_2, \ldots, x_n)\) 对 \(\overline{\mathbf{M}}\) 是正则的。通过对 \(n\) 作归纳推理，因此只需证明：若 \(\mathbf{M}\) 对 \(\mathfrak{x}\) -进拓扑是分离的且 \(\beta_{\mathbf{M}}^{\mathbf{x}}\) 是双射，则 \((x_1)_\mathbf{M}\) 是单射且 \(\beta_{\mathbf{M}}^{\mathbf{x}'}\) 是双射。现在， \(\beta_{\mathbf{M}}^{\mathbf{x}}\) 的双射性特别地蕴含了以 \(x_1\) 为比率、到 \(\bigoplus_{r} x^r \mathbf{M}/x^{r+1} \mathbf{M}\) 上的位似是单射的，从而 \(\text{Ker } (x_1)_\mathbf{M} \subset \bigcap_{i} x^i \mathbf{M}\) ，因此 \((x_1)_\mathbf{M}\) 是单射，如果 \(\mathfrak{x}\) -进拓扑在 \(\mathbf{M}\) 上是分离的。

因此我们归结为证明：如果 \((x_{1})_{\mathbf{M}}\) 是单射且 M 满足条件 \((\mathbf{v}')\) ，则 \(\overline{M}\) 满足条件 \((\mathbf{v}')\) （关于序列 \(x'\) ）。

根据假设，我们有一个正合列

\[
0 \longrightarrow \mathrm{M} \xrightarrow {(x _ {1}) _ {M}} \mathrm{M} \longrightarrow \overline {{\mathrm{M}}} \longrightarrow 0;
\]

令 \(\overline{\mathbf{A}}_{0} = \mathbf{A}_{0}/\mathbf{T}_{1}\mathbf{A}_{0},\quad\overline{\mathbf{t}} = \mathbf{t}\overline{\mathbf{A}}_{0}\) 。设 \(q:L\to M\) 为 \(\mathbf{A}_{0}\) -模 M 的一个自由分解；由于比率为 \(\mathbf{T}_{1}\) 的位似在 \(\mathbf{A}_{0}\) 中是单射，它在 L 中也是单射，并且我们有一个复形的正合列

\[
0 \longrightarrow \mathrm{L} \xrightarrow {(x _ {1}) _ {\mathrm{L}}} \mathrm{L} \longrightarrow \overline {{\mathrm{L}}} \longrightarrow 0
\]

其中 \(\overline{\mathbf{L}} = \mathbf{L} / x_1\mathbf{L}\) ，以及一个交换图

\[
\begin{array}{c} 0 \longrightarrow \mathrm{L} \stackrel {(x _ {1}) _ {\mathrm{L}}} {\longrightarrow} \mathrm{L} \longrightarrow \overline {{\mathrm{L}}} \longrightarrow 0 \\ \Big \downarrow^ {q} \quad \Big \downarrow^ {q} \quad \Big \downarrow^ {\overline {{q}}} \\ 0 \longrightarrow \mathrm{M} \stackrel {(x _ {1}) _ {\mathrm{M}}} {\longrightarrow} \mathrm{M} \longrightarrow \overline {{\mathrm{M}}} \longrightarrow 0. \end{array}
\]

由于 \(q\) 是一个拟同构， \(\overline{q}:\overline{\mathbf{L}}\to\overline{\mathbf{M}}\) 是 \(\overline{\mathbf{A}}_{0}\) -模 \(\overline{\mathbf{M}}\) 的一个自由分解（第 30 页，推论 1）。对每个 \(\overline{\mathbf{A}}_{0}\) -模 P，存在一个典范同构

\[
\mathbf {P} \otimes_ {\mathbf {A} _ {0}} \mathbf {L} \rightarrow \mathbf {P} \otimes_ {\overline {{{\mathbf {A}}}} _ {0}} \overline {{{\mathbf {L}}}},
\]

由此，通过取同调，得到一个同构

\[
\varphi_ {P}: \operatorname{Tor} _ {1} ^ {\mathbf {A} _ {0}} (\mathbf {P}, \mathbf {M}) \rightarrow \operatorname{Tor} _ {1} ^ {\overline {{\mathbf {A}}} _ {0}} (\mathbf {P}, \overline {{\mathbf {M}}}).
\]

若 \(u: \mathbf{P} \to \mathbf{P}'\) 是一个 \(\overline{\mathbf{A}}\) -同态，则

\[
\varphi_ {\mathrm{P} ^ {\prime}} \circ \operatorname{Tor} _ {1} ^ {\mathbf {A} _ {0}} (u, 1 _ {\mathrm{M}}) = \operatorname{Tor} _ {1} ^ {\overline {{\mathbf {A}}} _ {0}} (u, 1 _ {\mathrm{M}}) \circ \varphi_ {\mathrm{P}}.
\]

既然如此，假设条件 (v') 对 M 成立，下面我们来证明它对 \(\overline{M}\) 成立。设 \(r\) 为一个整数 \(\geqslant 1\) ；我们有一个行正合的交换图

\[
\begin{array}{c} 0 \longrightarrow \mathrm{A} _ {0} / t ^ {r} \xrightarrow {\mathrm{T} _ {1}} \mathrm{A} _ {0} / t ^ {r + 1} \longrightarrow \overline {{\mathrm{A}}} _ {0} / \overline {{t}} ^ {r + 1} \longrightarrow 0 \\ \downarrow^ {p _ {r - 1}} \quad \downarrow^ {p _ {r}} \quad \downarrow^ {\bar {p} _ {r}} \\ 0 \longrightarrow \mathrm{A} _ {0} / t ^ {r - 1} \xrightarrow {\mathrm{T} _ {1}} \mathrm{A} _ {0} / t ^ {r} \longrightarrow \overline {{\mathrm{A}}} _ {0} / \overline {{t}} ^ {r} \longrightarrow 0 \end{array}
\]

由此我们推出一个具有正合行的交换图

\[
\begin{array}{c c c c c} \operatorname{Tor} _ {1} ^ {\mathbf {A} _ {0}} (\mathbf {A} _ {0} / t ^ {r + 1}, \mathbf {M}) & \longrightarrow & \operatorname{Tor} _ {1} ^ {\mathbf {A} _ {0}} (\overline {{\mathbf {A}}} _ {0} / \overline {{t}} ^ {r + 1}, \mathbf {M}) & \longrightarrow & \mathbf {M} / x ^ {r}   \mathbf {M} \xrightarrow {x _ {1}} \mathbf {M} / x ^ {r + 1}   \mathbf {M} \\ \operatorname{Tor} _ {1} (p _ {r}, 1) & \downarrow & \operatorname{Tor} _ {1} (\bar {p} _ {r}, 1) & \downarrow & \downarrow \\ \operatorname{Tor} _ {1} ^ {\mathbf {A} _ {0}} (\mathbf {A} _ {0} / t ^ {r}, \mathbf {M}) & \longrightarrow & \operatorname{Tor} _ {1} ^ {\mathbf {A} _ {0}} (\overline {{\mathbf {A}}} _ {0} / \overline {{t}} ^ {r}, \mathbf {M}) & \longrightarrow & \mathbf {M} / x ^ {r - 1}   \mathbf {M} \xrightarrow {x _ {1}} \mathbf {M} / x ^ {r}   \mathbf {M}. \end{array}
\]

现在乘以 \(x_{1}\) 定义了从 \(\mathbf{M}/\mathfrak{x}^{r}\mathbf{M}\) 到 \(\mathbf{M}/\mathfrak{x}^{r+1}\mathbf{M}\) 的一个单射：通过对 \(r\) 归纳，这由正合列

\[
0 \rightarrow x ^ {r - 1} M / x ^ {r} M \rightarrow M / x ^ {r} M \rightarrow M / x ^ {r - 1} M \rightarrow 0
\]

以及比率为 \(x_{1}\) 的位似到 \(\bigoplus_{r\geqslant 0}(\mathfrak{x}^{r}\mathrm{M}/\mathfrak{x}^{r+1}\mathrm{M})\) 上的单射性立即得出。因此条件 (v) 蕴含同态 \(\operatorname{Tor}_{1}^{\mathbf{A}_{0}}(\overline{p}_{r},1_{\mathbf{M}})\) 对所有 \(r\geqslant 1\) 是满射的。通过与同构 \((\varphi_{\overline{\mathbf{A}}_{0}/\overline{\mathbf{t}}^{r+1}})^{-1}\) 和 \(\varphi_{\overline{\mathbf{A}}_{0}/\overline{\mathbf{t}}^{r}}\) 复合，我们推出同态

\[
\operatorname{Tor} _ {1} ^ {\overline {{\mathbf {A}}} _ {0}} \left(\bar {p} _ {r}, 1 _ {\mathbf {M}}\right): \operatorname{Tor} _ {1} ^ {\overline {{\mathbf {A}}} _ {0}} \left(\overline {{\mathbf {A}}} _ {0} / \bar {t} ^ {r + 1}, \mathbf {M}\right)\rightarrow \operatorname{Tor} _ {1} ^ {\overline {{\mathbf {A}}} _ {0}} \left(\overline {{\mathbf {A}}} _ {0} / \bar {t} ^ {r}, \mathbf {M}\right)
\]

对 \(r \geqslant 1\) 是满射的，从而得到 \(\overline{\mathbf{M}}\) 的条件 (v')。这就完成了定理的证明。

\subsection{与正则序列相伴的扩张类}

设 A 是一个环，M 是一个 A-模， \(x = (x_{1}, \ldots, x_{n})\) 是 A 中元素的序列。记 \(M_{i}\) 为 A-模 \(\mathbf{M}/(x_{1}\mathbf{M} + \cdots + x_{i-1}\mathbf{M})\) ，其中 \(i = 0, \ldots, n + 1\) ，使得 \(M_{0}\) 与 \(M_{1}\) 都与 M 等同，且 \(\mathbf{M}_{n+1} = \mathbf{M}/(\mathbf{x})\mathbf{M}\) 。设

\[
\overline {{x}} _ {i}: \mathbf {M} _ {i - 1} \to \mathbf {M} _ {i}, \qquad i = 1, \dots , n,
\]

表示由 \(M_{i-1}\) 上比率为 \(x_{i}\) 的位似与 \(M_{i-1}\) 到 \(M_{i}\) 上的典范投影复合而成的 A-同态。最后，我们记 \(p: M_{n} \to M/(x)M\) 为典范投影。图表

\[
0 \longrightarrow \mathrm{M} \xrightarrow {\overline {{{x}}} _ {1}} \mathrm{M} _ {1} \xrightarrow {\overline {{{x}}} _ {2}} \mathrm{M} _ {2} \longrightarrow \dots \xrightarrow {\overline {{{x}}} _ {n}} \mathrm{M} _ {n} \xrightarrow {p} \mathrm{M} / (\boldsymbol {x}) \mathrm{M} \longrightarrow 0\tag{19}
\]

是一个正合列当且仅当序列 x 是 M-正则的。此后假设序列 x 对 M 是正则的。与正合列 (19) 相伴的元素 \(\theta_{x} \in \operatorname{Ext}_{A}^{n}(\mathbf{M}/(\mathbf{x})\mathbf{M}, \mathbf{M})\) 也称为与 M-正则序列 x 相伴。

设 i 为整数， \(1 \leqslant i \leqslant n\) 。注意，序列 (19) 可以分解为两个正合列

\begin{enumerate}
\def\labelenumi{(\arabic{enumi})}
\setcounter{enumi}{19}
\tightlist
\item
\end{enumerate}

\[
0 \longrightarrow \mathrm{M} \xrightarrow {\overline {{x}} _ {1}} \mathrm{M} _ {1} \xrightarrow {\overline {{x}} _ {2}} \mathrm{M} _ {2} \longrightarrow \dots \xrightarrow {\overline {{x}} _ {i}} \mathrm{M} _ {i} \longrightarrow \mathrm{M} / (x _ {1} \mathrm{M} + \dots + x _ {i} \mathrm{M}) \longrightarrow 0\tag{21}
\]

\[
0 \longrightarrow \mathrm{M} / (x _ {1} \mathrm{M} + \dots + x _ {i} \mathrm{M}) \xrightarrow {\overline {{x}} _ {i + 1}} \mathrm{M} _ {i + 1} \longrightarrow \dots
\]

\[
\longrightarrow \mathrm{M} _ {n} \stackrel {p} {\longrightarrow} \mathrm{M} / (x) \mathrm{M} \longrightarrow 0
\]

它们正是与序列 \((x_{1}, \ldots, x_{i})\) （它对 M 是正则的）以及与序列 \((x_{i+1}, \ldots, x_{n})\) （它对 \(\mathrm{M}/(x_{1} \mathrm{~M} + \cdots + x_{i} \mathrm{~M})\) 是正则的）相伴的正合列。记

\[
\theta_ {(x _ {1}, \dots , x _ {i})} \in \operatorname{Ext} _ {\mathbf {A}} ^ {i} (\mathbf {M} / (x _ {1} \mathbf {M} + \dots + x _ {i} \mathbf {M}), \mathbf {M})
\]

\[
\theta_ {(x _ {i + 1}, \dots , x _ {n})} \in \operatorname{Ext} _ {\mathbf {A}} ^ {n - i} (\mathbf {M} / (\boldsymbol {x}) \mathbf {M}, \mathbf {M} / (x _ {1} \mathbf {M} + \dots + x _ {i} \mathbf {M}))
\]

为与 (20) 和 (21) 相伴的扩张类，则根据 X，第 118 页，命题 3，我们有

\[
\theta_ {(x _ {1}, \dots , x _ {n})} = \theta_ {(x _ {1}, \dots , x _ {i})} \circ \theta_ {(x _ {i + 1}, \dots , x _ {n})}.\tag{22}
\]

此外，根据命题 5（X，第 157 页），科斯居尔复形 K.(x, M) 在次数 \(\neq n\) 时零调，由此得到一个正合列

\[
0 \longrightarrow \mathrm{M} \xrightarrow {\partial^ {0}} \mathbf {K} ^ {1} (\boldsymbol {x}, \mathrm{M}) \xrightarrow {\partial^ {1}} \mathbf {K} ^ {2} (\boldsymbol {x}, \mathrm{M}) \xrightarrow {\partial^ {2}} \dots \longrightarrow \mathbf {K} ^ {n} (\boldsymbol {x}, \mathrm{M}) \xrightarrow {q} \mathrm{M} / (\boldsymbol {x}) \mathrm{M} \longrightarrow 0\tag{23}
\]

其中我们已把 \(\mathbf{K}^{0}(\boldsymbol{x},\mathbf{M})\) 与 \(\mathbf{M}\) 等同，且其中 \(q\) 把每个元素 \(\boldsymbol{m}\in\mathbf{K}^{n}(\boldsymbol{x},\mathbf{M})\) 映到 \(\mathbf{M}/(\boldsymbol{x})\mathbf{M}\) 中 \(\boldsymbol{m}(1,2,\ldots,n)\in\mathbf{M}\) 所在的类。

命题 8. --- 假设序列 x 对 M 是正则的。Ext \(_{A}^{n}\) (M/(x) M, M) 中与正合列 (23) 相伴的元素是 (-1) \(^{n(n+1)/2}\) . θ \(_{x}\) 。

对 \(i = 0,1,\dots,n\) ，让我们定义一个 A-线性映射

\[
a ^ {i}: \mathbf {K} ^ {i} (\boldsymbol {x}, \mathrm{M}) \rightarrow \mathrm{M} _ {i} = \mathrm{M} / (x _ {1} \mathrm{M} + \dots + x _ {i - 1} \mathrm{M})
\]

如下：如果 \(\boldsymbol{m} \in \mathbf{K}^i(\boldsymbol{x}, \mathbf{M})\) ，则 \(a^i(\boldsymbol{m})\) 是 \(\mathbf{M}_i\) 中元素 \(\boldsymbol{m}(1, 2, \ldots, i)\) （属于 M）的类。显然， \(a^0\) 是 M 的恒等映射，并且 \(p \circ a^n = q\) 。此外， \(a^{i+1} \circ \partial^i(\boldsymbol{m})\) 是下述元素在 \(\mathbf{M}_{i+1}\) 中的像

\[
\sum_ {k = 1} ^ {i + 1} (- 1) ^ {k + 1} x _ {k} \boldsymbol {m} (1, 2, \dots , k - 1, k + 1, \dots , i + 1).
\]

由于 \(x_{k}\) 零化 \(\mathbf{M}_{i+1}\) 对 \(k=1,\ldots,i,a^{i+1}\circ\partial^{i}(\boldsymbol{m})\) 也是下述元素的像

\[
(- 1) ^ {i} x _ {i + 1} \boldsymbol {m} (1, 2, \dots , i),
\]

因此

\[
a ^ {i + 1} \circ \partial^ {i} = (- 1) ^ {i} \overline {{{x}}} _ {i + 1} \circ a ^ {i}.
\]

根据 X，第 120 页，推论 1 和 2，与 (23) 相伴的 Ext \(_{A}^{n}\) (M/(x) M, M) 的元素等于 \(\prod_{i=1}^{n} (-1)^{i} \cdot \theta_{x}\) ，由此得该断言。

推论。 — 进一步假设诸模 \(\mathbf{M}/(x_{1}\mathbf{M}+\cdots+x_{i-1}\mathbf{M})\) 对 (x)-进拓扑是分离的，并设 \((a_{ij})\in\mathbf{GL}_{n}(\mathbf{A})\) 。令

\[
y _ {i} = \sum_ {j} a _ {i j} x _ {j} \text { and } \mathbf {y} = (y _ {1}, \dots , y _ {n}).
\]

则序列 y 对 M 是正则的，并且有 \(\theta_{y} = \det (a_{ij})^{-1}\theta_{x}\) 。

事实上，序列 y 对 M 是正则的，由命题 6（X，第 158 页）和定理 1（X，第 160 页）可得；最后一个断言由命题 8 以及 X 第 119 页的命题 4 得出。

命题 9. --- 假设序列 x 对 M 是正则的。若 N 是满足 (x) N = 0 的 A-模，则我们有 Ext \(_{A}^{i}\) (N, M) = 0 对 i \textless{} n，并且映射 α → θ \(_{x}\) ∘ α 从 Hom \(_{A}\) (N, M/(x) M) 到 Ext \(_{A}^{n}\) (N, M) （它也是与 (19) 相伴的迭代连接同态，参见 X，第 127 页，推论 3）是双射的。

仍需证明同态 \(\psi^{i}:\alpha\mapsto\theta_{x}\circ\alpha\) 从 \(\operatorname{Ext}_{\mathbf{A}}^{i-n}(\mathbf{N},\mathbf{M}/(\mathbf{x})\mathbf{M})\) 到 \(\operatorname{Ext}_{\mathbf{A}}^{i}(\mathbf{N},\mathbf{M})\) 对 \(i\leqslant n\) 是双射的。让我们对 \(n\) 作归纳推理，该断言对 \(n=0\) 是平凡的。令 \(\mathbf{M}_{1}=\mathbf{M}/x_{1}\mathbf{M}\) , \(\mathbf{x}^{\prime}=(x_{2},...,x_{n})\) ，使得 \(\mathbf{x}^{\prime}\) 对 \(\mathbf{M}_{1}\) 是正则的。由归纳假设，同态

\[
\overline {{{{\psi}}}} ^ {i - 1}: \alpha \mapsto \theta_ {\boldsymbol {x} ^ {\prime}} \circ \alpha
\]

从 \(\operatorname{Ext}_{\mathbf{A}}^{i-n}(\mathbf{N}, \mathbf{M}/(\mathbf{x})\mathbf{M})\) 到 \(\operatorname{Ext}_{\mathbf{A}}^{i-1}(\mathbf{N}, \mathbf{M}_1)\) 对 \(i < n\) 是双射的。此外，考虑正合列

\[
0 \longrightarrow \mathrm{M} \xrightarrow {(x _ {1}) _ {\mathrm{M}}} \mathrm{M} \longrightarrow \mathrm{M} _ {1} \longrightarrow 0;
\]

与之相伴的连接同态 \(\operatorname{Ext}_{\mathbf{A}}^{i}(\mathbf{N}, \mathbf{M}_{1}) \to \operatorname{Ext}_{\mathbf{A}}^{i+1}(\mathbf{N}, \mathbf{M})\) 为 \(\varphi^{i}: \beta \mapsto \theta_{x_{1}} \circ \beta\) （X，第 125 页，命题 5）；

由于我们有 \(\operatorname{Ext}^i(1_{\mathbf{N}}, (x_1)_{\mathbf{M}}) = \operatorname{Ext}^i((x_1)_{\mathbf{N}}, 1_{\mathbf{M}}) = 0\) ，我们推出正合列

\[
0 \longrightarrow \operatorname{Ext} _ {\mathrm{A}} ^ {i} (\mathrm{N}, \mathrm{M}) \longrightarrow \operatorname{Ext} _ {\mathrm{A}} ^ {i} (\mathrm{N}, \mathrm{M} _ {1}) \xrightarrow {\varphi^ {i}} \operatorname{Ext} _ {\mathrm{A}} ^ {i + 1} (\mathrm{N}, \mathrm{M}) \longrightarrow 0.
\]

由于 \(\operatorname{Ext}_{\mathbf{A}}^{i}(\mathbf{N}, \mathbf{M}_{1}) = 0\) 对 \(i < n - 1\) ，我们推出 \(\operatorname{Ext}_{\mathbf{A}}^{i + 1}(\mathbf{N}, \mathbf{M}) = 0\) 对 \(i < n - 1\) ，即 \(i + 1 < n\) ；由此可知 \(\varphi^{i}\) 对 \(i < n\) 是双射的。由于 \(\psi^{i}(\alpha) = \theta_{x} \circ \alpha = \theta_{x_{1}} \circ \theta_{x'} \circ \alpha = \varphi^{i - 1} \circ \overline{\psi}^{i}(\alpha)\) 对 \(\alpha \in \operatorname{Ext}_{\mathbf{A}}^{i - n}(\mathbf{N}, \mathbf{M}/(\mathbf{x})\mathbf{M})\) ， \(\psi^{i}\) 对 \(i \leqslant n\) 确实是双射的。

\backmatter
\specialchapter{习题}
